\motto{Politiques et scientifiques ont le sens des réalités, mais ce ne sont pas les mêmes. Il en résulte --- et ce sera là un principe que le général de Gaulle fera sien que l'activité de recherche ne peut être évaluée, quant à sa qualité propre, que par des hommes qui la pratiquent eux-mêmes.

--- Pierre Lelong\footnote{Pierre Lelong (1912--2011) was the husband of another famous mathematician Jacqueline Ferrand. During their marriage (1947--1977), the latter published under the name of Jacqueline Lelong-Ferrand.}, 1999}
\chapter{The envelope operators}
\label{chap:envelopes}

This chapter develops the envelope formalism that will be used throughout the rest of the book. There are two parallel constructions. The first is the $P$-envelope, which records the singularity type relevant to non-pluripolar Monge--Amp\`ere mass. The second is the $\mathcal I$-envelope, which records the singularity type detected by multiplier ideal sheaves. The distinction between these two envelopes is one of the main organizing principles of the book.

The starting point is the rooftop operator. Given two singular potentials, the rooftop construction produces the largest potential lying below both, and it provides a flexible way to compare singularities. From this we define the $P$-envelope $P_\theta[\varphi]$. Its fixed points are the model potentials. Model potentials are the natural representatives of $P$-singularity types: they preserve non-pluripolar mass and satisfy strong contact-set properties. The first part of the chapter establishes the basic algebra of these envelopes, their behavior under monotone limits, and the domination and contact principles needed later.

We then introduce the Monge--Amp\`ere energy relative to a model potential. This leads to the finite-energy classes $\mathcal E^1(X,\theta;\phi)$ and to stability results for rooftop envelopes. These results will be used in the metric theory of singularities and in the study of geodesic rays.

The second part of the chapter develops the $\mathcal I$-envelope $P_\theta[\varphi]_{\mathcal I}$. This envelope is defined by imposing multiplier ideal conditions rather than ordinary singularity-type conditions. It is therefore sensitive to valuative data in a way that the ordinary $P$-envelope is not. We prove its basic functorial properties, its independence of the chosen representative of the cohomology class, its bimeromorphic behavior, and its compatibility with analytic singularities.

The comparison between $P_\theta[\varphi]$ and $P_\theta[\varphi]_{\mathcal I}$ is delicate. They agree in important cases, for instance for analytic singularities, but they differ in general. Later chapters use this comparison to define $\mathcal I$-good singularities, trace operators, partial Okounkov bodies and non-Archimedean potentials. Thus the technical results of this chapter should be regarded as the basic calculus of singularity types used in the sequel.

Several results about $P$-envelopes are developed by Darvas--Di Nezza--Lu in various papers and summarized in their survey \cite{DDNLsurv}. The reader is invited to read their survey for more details and further results.

\section{The \texorpdfstring{$P$}{P}-envelope}\label{sec:Penv}
In this section, $X$ will denote a connected compact Kähler manifold of dimension $n$.


\subsection{Rooftop operator and the definition of the \texorpdfstring{$P$}{P}-envelope}
\label{subsec:rooftop_Penv}

Fix a smooth closed real $(1,1)$-form $\theta$ on $X$.


\begin{definition}\label{def:rooftop}
    Given $\varphi,\psi\in \PSH(X,\theta)$, we define their \emph{rooftop operator}\index{rooftop operator} as follows:
    \begin{equation}\label{eq:rooftopdef1}
    \varphi\land\psi=\sup\left\{\eta\in \PSH(X,\theta):\eta\leq \varphi,\eta\leq \psi\right\}.
    \end{equation}
    \index{$\varphi\land \psi$}
    For simplicity of notation, we extend the definition to the case where $\varphi$ or $\psi$ is constantly $-\infty$, in which case we simply set
    \[
    \varphi\land \psi=-\infty.
    \]
\end{definition}
When we want to be more specific, we may also write $\varphi\land_{\theta}\psi$. Note that it is possible that the set in \eqref{eq:rooftopdef1} is empty, in which case $\varphi\land \psi=-\infty$.


\begin{proposition}\label{prop:land_well_def}
    The operator $\land$ is a well-defined commutative, associative binary operator
    \[
    \bigl(\PSH(X,\theta)\cup\{-\infty\}\bigr)\times\bigl(\PSH(X,\theta)\cup\{-\infty\}\bigr)\rightarrow \PSH(X,\theta)\cup\{-\infty\}.
    \]
\end{proposition}
\begin{proof}
    We first show that the map is well-defined. For this purpose, take $\varphi,\psi\in \PSH(X,\theta)$. When the set in \eqref{eq:rooftopdef1} is empty, there is nothing to prove. So let us assume that the set is not empty.

    Define
    \[
    \gamma=\uscsup \left\{\eta\in \PSH(X,\theta):\eta\leq \varphi,\eta\leq \psi\right\}.
    \]
    Then by \cref{prop:pshfunction_closedseq}, we find that $\gamma\in \PSH(X,\theta)$ and hence $\gamma$ is a candidate for the supremum in \eqref{eq:rooftopdef1}. Therefore, $\gamma\leq \varphi\land \psi$. The reverse inequality follows from the definition of $\gamma$, so
    \[
    \varphi\land \psi=\gamma\in \PSH(X,\theta).
    \]
    The commutativity and the associativity of $\land$ follow directly from the defining maximality property.
\end{proof}
Note that from the proof, we have
\[
\varphi\land \psi=P_{\theta}\Bigl( \min\{\varphi,\psi\} \Bigr).
\]
See \cref{def:Pthetaf} for the notation $P_{\theta}(\bullet)$.


\begin{lemma}\label{lma:rooftopMA}
    Let $\varphi,\psi\in \PSH(X,\theta)$. Assume that $\varphi\land\psi\in \PSH(X,\theta)$. Then
    \begin{equation}\label{eq:rooftopMAineq1}
    \theta_{\varphi\land \psi}^n\leq \mathds{1}_{\{\varphi\land\psi=\varphi\}}\, \theta_{\varphi}^n+\mathds{1}_{\{\varphi\land\psi=\psi\}} \,\theta_{\psi}^n.
    \end{equation}
\end{lemma}
In particular, $\theta_{\varphi\land \psi}^n$ is carried by the union $\{\varphi\land \psi=\varphi\}\cup \{\varphi\land \psi=\psi\}$.

\begin{proof}
We first observe that as a consequence of \cref{lma:lorMAineqe}, we have
\[
\mathds{1}_{\{\varphi\land \psi=\varphi\}} \,\theta_{\varphi\land \psi}^n\leq \mathds{1}_{\{\varphi\land \psi=\varphi\}} \,\theta_{\varphi}^n,\quad \mathds{1}_{\{\varphi\land \psi=\psi\}} \,\theta_{\varphi\land \psi}^n\leq \mathds{1}_{\{\varphi\land \psi=\psi\}} \,\theta_{\psi}^n.
\]
Applying \cref{thm:concentration} to $\min\{\varphi,\psi\}$, we conclude that
\[
\theta_{\varphi\land \psi}^n\leq \mathds{1}_{\{\varphi\land \psi=\varphi\}} \,\theta_{\varphi\land \psi}^n+\mathds{1}_{\{\varphi\land \psi=\psi\}}\,\theta_{\varphi\land \psi}^n\leq \mathds{1}_{\{\varphi\land\psi=\varphi\}} \,\theta_{\varphi}^n+\mathds{1}_{\{\varphi\land\psi=\psi\}} \,\theta_{\psi}^n,
\]
and \eqref{eq:rooftopMAineq1} is established.
\end{proof}



We recall that the relations $\preceq$ and $\sim$ are introduced in \cref{def:parorder}.
\begin{definition}\label{def:Penv}
    Given $\varphi\in \PSH(X,\theta)$, we define its \emph{$P$-envelope}\index{envelope!P-envelope@$P$-envelope} as follows:
    \begin{equation}\label{eq:Pthetavarphi}
    P_{\theta}[\varphi]\coloneqq \uscsup \left\{ \psi\in \PSH(X,\theta): \psi\leq 0,\psi\preceq \varphi \right\}.
    \end{equation}
    \index{$P_{\theta}[\varphi]$}
\end{definition}
Note that $\varphi-\sup_X \varphi$ is always a candidate for the supremum. By \cref{prop:pshfunction_closedseq}, we have $P_{\theta}[\varphi]\in \PSH(X,\theta)$ and $P_{\theta}[\varphi]\leq 0$.
Moreover, the definition can be equivalently described as
\begin{equation}\label{eq:Penvsups}
    P_{\theta}[\varphi]=\uscsup[C\in \mathbb{Z}_{>0}](\varphi+C)\land V_{\theta}.
\end{equation}
Recall that $V_{\theta}$ is introduced in \eqref{eq:Vtheta}. For any $C\in \mathbb{R}$, we have $(\varphi+C)\land V_{\theta}\in \PSH(X,\theta)$ and
\[
(\varphi+C)\land V_{\theta}\sim \varphi.
\]
In other words, in \eqref{eq:Pthetavarphi}, we may replace the condition $\psi\preceq \varphi$ by $\psi\sim \varphi$. Also note that as a consequence,
\begin{equation}\label{eq:varphipreceqP}
    \varphi\preceq P_{\theta}[\varphi].
\end{equation}



The idea lying behind the definition of $P_{\theta}[\varphi]$ is that we choose the least singular element out of all potentials with the same singularity type as $\varphi$. As we shall see in \cref{ex:loglogsing} below, $P_{\theta}[\varphi]$ does not necessarily have the same singularity type as $\varphi$. This forces us to define a rougher equivalence relation in \cref{def:Pmoresing}.

The envelope depends on the choice of $\theta$, but the dependence is easy to understand:
\begin{proposition}\label{prop:Penvindeptheta}
    Let $\theta'=\theta+\ddc g$ for some $g\in C^{\infty}(X)$. Then for any $\varphi\in \PSH(X,\theta)$, we have $\varphi'\coloneqq \varphi-g\in \PSH(X,\theta')$ and
    \[
        P_{\theta}[\varphi]\sim P_{\theta'}[\varphi'].
    \]
\end{proposition}
\begin{proof}
    By symmetry, it suffices to show that
    \[
        P_{\theta}[\varphi]\preceq P_{\theta'}[\varphi'].
    \]
    We may assume that $g\geq 0$ (after replacing $g$ by $g+C$ with $C\geq -\inf_X g$, which leaves $\theta'$ unchanged and shifts $\varphi'$ by an additive constant).
    Then for any $\psi\in \PSH(X,\theta)$ with $\psi\preceq \varphi$ and $\psi\leq 0$, we set $\psi'\coloneqq \psi-g\in \PSH(X,\theta')$. Then $\psi'\preceq \varphi'$ and $\psi'\leq 0$, so $\psi'\leq P_{\theta'}[\varphi']$. Since $\psi$ is arbitrary, it follows that
    \[
        P_{\theta}[\varphi]-\sup_X g\leq P_{\theta}[\varphi]-g\leq P_{\theta'}[\varphi'],
    \]
    which gives $P_{\theta}[\varphi]\preceq P_{\theta'}[\varphi']$.
\end{proof}

The $P$-envelope preserves the non-pluripolar masses:
\begin{proposition}\label{prop:Ppresmass}
Suppose that $\theta_1,\ldots,\theta_n$ are smooth closed real $(1,1)$-forms on $X$. Let $\varphi_i\in \PSH(X,\theta_i)$ for each $i=1,\ldots,n$. Then
\begin{equation}\label{eq:Penvpremass}
\int_X \theta_{1,P_{\theta_1}[\varphi_1]}\wedge \cdots \wedge \theta_{n,P_{\theta_n}[\varphi_n]}=\int_X \theta_{1,\varphi_1}\wedge \cdots\wedge \theta_{n,\varphi_n}.
\end{equation}
\end{proposition}
This proposition together with \cref{thm:mono} will be combined to a common generalization in \cref{prop:mono2} after introducing the $P$-preorder.
\begin{proof}
For each $C\in \mathbb{Z}_{>0}$ and each $i=1,\ldots,n$, we have
\[
(\varphi_i+C)\land V_{\theta_i}\sim \varphi_i.
\]
It follows from the monotonicity theorem \cref{thm:mono} that
\[
\int_X \theta_{1,(\varphi_1+C)\land V_{\theta_1}}\wedge \cdots \wedge \theta_{n,(\varphi_n+C)\land V_{\theta_n}}=\int_X \theta_{1,\varphi_1}\wedge \cdots\wedge \theta_{n,\varphi_n}.
\]
So \eqref{eq:Penvpremass} follows from \eqref{eq:Penvsups} and \cref{cor:incseqnppcont}.
\end{proof}

Conversely, \cref{prop:Ppresmass} characterizes the $P$-envelope, as we will see in a moment in \cref{thm:Pvarphidiffdef}. We need some preparation for the proof. We first show that the $P$-envelope can be regarded as a concentration of the non-pluripolar mass:
\begin{theorem}\label{thm:Pvarphisupport}
    Let $\varphi\in \PSH(X,\theta)$. Then
    \begin{equation}\label{eq:thetaPthetaphin}
        \theta_{P_{\theta}[\varphi]}^n\leq \mathds{1}_{\{P_{\theta}[\varphi]=0\}}\,\theta^n.
    \end{equation}
\end{theorem}
In fact, equality holds as we will see in \cref{prop:DNTenvelope2} below.
\begin{proof}
    Thanks to \cref{lma:rooftopMA}, for each $C>0$, we have
    \[
    \begin{aligned}
    \theta_{(\varphi+C)\land V_{\theta}}^n\leq & \mathds{1}_{\{(\varphi+C)\land V_{\theta}=\varphi+C\} }\, \theta_{\varphi}^n+\mathds{1}_{\{(\varphi+C)\land V_{\theta}=V_{\theta}\} }\, \theta_{V_{\theta}}^n\\
    \leq & \mathds{1}_{\{\varphi+C\leq V_{\theta}\} }\, \theta_{\varphi}^n+\mathds{1}_{\{P_{\theta}[\varphi]=V_{\theta}\} } \,\theta_{V_{\theta}}^n.
    \end{aligned}
    \]
    We wish to let $C\to\infty$.
    The dominated convergence theorem and the semicontinuity theorem \cref{cor:incseqnppcont} give, as $C\to\infty$,
    \[
    \mathds{1}_{\{\varphi+C\leq V_{\theta}\} } \theta_{\varphi}^n\rightharpoonup 0,\quad
    \theta_{(\varphi+C)\land V_{\theta}}^n\rightharpoonup \theta_{P_{\theta}[\varphi]}^n.
    \]
    So we conclude that
    \[
    \theta_{P_{\theta}[\varphi]}^n\leq \mathds{1}_{\{P_{\theta}[\varphi]=V_{\theta}\}} \,\theta_{V_{\theta}}^n.
    \]
    Taking \cref{cor:DNTenvelope1} into consideration, we conclude \eqref{eq:thetaPthetaphin}.
\end{proof}

Using essentially the same proof, we arrive at the following conclusion:
\begin{corollary}\label{cor:Pvarphisupport}
    Let $\varphi,\psi\in \PSH(X,\theta)$. Let
    \[
    \eta\coloneqq \uscsup[C>0] (\varphi+C)\land \psi.\footnotemark
    \]
    \footnotetext{Using the terminology to be introduced in \cref{def:Pthetavarphif}, we can also write $\eta=P_{\theta}[\varphi](\psi)$.}
    Assume that $\eta\not\equiv -\infty$ (for example when $\varphi\preceq \psi$).
    Then
    \[
    \theta_{\eta}^n\leq \mathds{1}_{\{\eta=\psi\}}\,\theta_{\psi}^n.
    \]
\end{corollary}
\begin{proof}
    Thanks to \cref{lma:rooftopMA}, for each $C>0$, we have
    \[
    \begin{aligned}
    \theta_{(\varphi+C)\land \psi}^n\leq & \mathds{1}_{\{(\varphi+C)\land \psi=\varphi+C\} }\, \theta_{\varphi}^n+\mathds{1}_{\{(\varphi+C)\land \psi=\psi\} }\, \theta_{\psi}^n\\
    \leq & \mathds{1}_{\{\varphi+C\leq \psi\} }\, \theta_{\varphi}^n+\mathds{1}_{\{\eta=\psi\} } \,\theta_{\psi}^n.
    \end{aligned}
    \]
    We wish to let $C\to\infty$.
    The dominated convergence theorem and the semicontinuity theorem \cref{cor:incseqnppcont} give, as $C\to\infty$,
    \[
    \mathds{1}_{\{\varphi+C\leq \psi\} } \theta_{\varphi}^n\rightharpoonup 0,\quad
    \theta_{(\varphi+C)\land \psi}^n\rightharpoonup \theta_{\eta}^n.
    \]
    So we conclude that
    \[
    \theta_{\eta}^n\leq \mathds{1}_{\{\eta=\psi\}} \,\theta_{\psi}^n.
    \]
    This completes the proof.
\end{proof}


\begin{theorem}\label{thm:Pvarphidiffdef}
    Assume that $\varphi\in \PSH(X,\theta)_{>0}$. Then
    \begin{equation}\label{eq:Penvdef}
    P_{\theta}[\varphi]=\sup\left\{\psi\in \PSH(X,\theta): \psi\leq 0,\varphi\preceq \psi, \int_X \theta_{\varphi}^n=\int_X \theta_{\psi}^n\right\}.
    \end{equation}
    In particular, in this case,\footnote{Pang--Sun--Wang--Zhou \cite{PSWZ26} recently proved that, \eqref{eq:Penvprojop} continues to hold even when $\varphi$ has vanishing mass.}
    \begin{equation}\label{eq:Penvprojop}
    P_{\theta}\Bigl[P_{\theta}[\varphi]\Bigr]=P_{\theta}[\varphi].
    \end{equation}
\end{theorem}
Note that in \eqref{eq:Penvdef} and \eqref{eq:Pthetavarphi}, the auxiliary function $\psi$ lies on different sides of $\varphi$.
\begin{proof}
     Let $\psi$ be a candidate of the right-hand side of \eqref{eq:Penvdef}. It follows from \cref{thm:Pvarphisupport} that
     \[
     \int_{\{P_{\theta}[\varphi]<\psi \}} \theta_{P_{\theta}[\varphi]}^n\leq \int_{\{P_{\theta}[\varphi]<\psi \}\cap \{P_{\theta}[\varphi]=0\} } \theta^n=0.
     \]
     On the other hand, we have
    \[
    P_{\theta}[\varphi]\preceq P_{\theta}[\psi],\quad \psi\preceq P_{\theta}[\psi]
    \]
    and all these potentials have the same mass due to \cref{prop:Ppresmass}. So the domination principle \cref{thm:dom_prin_1} is applicable and gives
    \[
    P_{\theta}[\varphi]\geq \psi.
    \]
    Hence we get the $\geq$ direction in \eqref{eq:Penvdef}.

    On the other hand, $P_{\theta}[\varphi]$ itself is a candidate of the right-hand side of \eqref{eq:Penvdef}, as a consequence of \cref{prop:Ppresmass}. Therefore, \eqref{eq:Penvdef} follows.

    As for \eqref{eq:Penvprojop}, the $\geq$ direction is trivial. For the other direction, let $\psi\in \PSH(X,\theta)$ be a potential satisfying
    \[
    \psi\leq 0,\quad P_{\theta}[\varphi]\preceq \psi,\quad \int_X \theta_{P_{\theta}[\varphi]}^n=\int_X \theta_{\psi}^n.
    \]
    Then it follows from \cref{prop:Ppresmass} that
    \[
    \psi\leq 0,\quad \varphi\preceq \psi,\quad \int_X \theta_{\varphi}^n=\int_X \theta_{\psi}^n.
    \]
    In view of \eqref{eq:Penvdef}, we conclude the $\leq$ direction of \eqref{eq:Penvprojop}.
\end{proof}

\begin{definition}\label{def:modelpot}
    If $\varphi=P_{\theta}[\varphi]$ and $\int_X \theta_{\varphi}^n>0$, we say $\varphi$ is a \emph{model potential}.\index{potential!model potential}
\end{definition}
Recall that the notion of model potentials depends heavily on the choice of $\theta$. When there is a risk of confusion, we also say $\varphi$ is a model potential in $\PSH(X,\theta)$.

\begin{remark}\label{rmk:envelopes:L267}
\cref{def:modelpot} is different from the common definition in the literature: We impose the extra condition $\int_X \theta_{\varphi}^n>0$. The author believes that this is the only case where this notion is natural. We sometimes emphasize this point by saying $\varphi\in \PSH(X,\theta)_{>0}$ is a model potential.
\end{remark}

There are plenty of model potentials:
\begin{corollary}\label{cor:Psendspotentialtomodel}
    Let $\varphi\in \PSH(X,\theta)_{>0}$. Then $P_{\theta}[\varphi]$ is a model potential in $\PSH(X,\theta)$. Moreover,
    \[
        \int_X \theta_{P_{\theta}[\varphi]}^n=\int_X \theta_{\varphi}^n.
    \]
\end{corollary}
\begin{proof}
    This follows immediately from \cref{thm:Pvarphidiffdef} and \cref{prop:Ppresmass}.
\end{proof}

As we have already seen in the proof of \cref{lma:pathoenvelope}, we have the following interesting property:
\begin{proposition}\label{prop:model_sup_prop}
    Let $\varphi\in \PSH(X,\theta)$. Consider a model potential $\phi\in \PSH(X,\theta)_{>0}$. Assume that $\varphi\preceq \phi$. Then
    \[
    \sup_{\{\phi\neq -\infty\}} (\varphi-\phi)=\sup_X \varphi.
    \]
\end{proposition}
In particular, for any $\varphi\in \PSH(X,\theta)$, $\varphi\preceq \phi$ and $\varphi\leq 0$ implies $\varphi\leq \phi$.

\begin{proof}
    Observe that $\phi\leq 0$, so on the set $\{\phi\neq -\infty\}$, we have
    \[
    \varphi-\phi\geq \varphi.
    \]
    So the $\geq$ direction follows from \cref{cor:diffsupinfindeppluripolar}.
    Conversely, by assumption, we can find a constant $C>0$ so that
    \[
    \varphi-\sup_X \varphi\leq \min\{\phi+C,0 \}.
    \]
    It follows that $\varphi-\sup_X \varphi\leq P_{\theta}[\phi]=\phi$. Therefore, the reverse inequality follows.
\end{proof}



\begin{example}\label{ex:loglogsing}
    We continue our favorite example \cref{ex:LPP1}. Let $X=\mathbb{P}^1$ and $\omega$ be the Fubini--Study metric. We define $\varphi\in \PSH(X,\omega)$ as follows: For $z\in \mathbb{C}$, we let
    \[
    \varphi(z)=
    \begin{cases}
        -\log \left(|z|^2+1\right)+\Bigl(-\log \left(-\log |z|^2\right)\Bigr) \lor \Bigl(2+\log |z|^2\Bigr),& \textup{if }|z|<1/\sqrt{2},\\
        2+\log \frac{|z|^2}{|z|^2+1},& \textup{otherwise},
    \end{cases}
    \]
    while $\varphi(\infty)=2$. Note that $\varphi$ is $\omega$-psh on $X=\mathbb{P}^1$ by the gluing lemma \cref{lma:gl}.
The singularity of $\varphi$ only occurs at $z=0$, close to which, $\varphi$ is approximately $-\log \left(-\log |z|^2\right)$. This type of singularity is therefore called the \emph{log-log type singularity}.

    We claim that
    \begin{equation}\label{eq:Pomegavarphi0temp1}
        P_{\omega}[\varphi]=0.
    \end{equation}
    In particular, we find that $\varphi$ and $P_{\omega}[\varphi]$ have different singularity types.

    By \cref{thm:Pvarphidiffdef}, in order to verify \eqref{eq:Pomegavarphi0temp1}, it suffices to verify that
    \begin{equation}\label{eq:Pomegavarphi0temp2}
    \int_X \omega_{\varphi}^1=1.
    \end{equation}
    Here $\omega_{\varphi}^1$ is taken in the non-pluripolar sense. Since $\{0,\infty\}\subseteq \mathbb{P}^1$ is pluripolar, this reduces to showing that
    \[
    \int_{\mathbb{C}^*} \ddc \psi=\frac{1}{4\pi}\int_{\mathbb{C}^*} (\Delta \psi)\,\mathrm{d}\mu=1,
    \]
    where $\psi(z)=\varphi(z)+\log (|z|^2+1)$ and $\mu$ is the standard Lebesgue measure on $\mathbb{C}$.

    Note that the Laplacian vanishes outside $\overline{B(0,0.7)}$ since $\psi(z)=2+\log |z|^2$ there, which is harmonic. Therefore,
    \[
    \int_{\mathbb{C}^*} \ddc \psi=\frac{1}{4\pi}\int_{|z|<1/\sqrt{2}} (\Delta \psi)(z)\,\mathrm{d}\mu.
    \]
    It is an elementary exercise to see that the right-hand side is exactly equal to $1$. If you are familiar with toric geometry, this is more or less trivial since
    \[
    \nabla_r \Bigl( \left(-\log (-r)\right)\lor (2+r) \Bigr)\bigl((-\infty,-\log 2)\bigr)= (0,1].
    \]
    Otherwise, just try to evaluate the integral using Green's identities. Therefore, \eqref{eq:Pomegavarphi0temp2} is proved and our assertion \eqref{eq:Pomegavarphi0temp1} follows.
\end{example}


Next we give a criterion on when the rooftop operator is not identically $-\infty$.
\begin{proposition}\label{prop:landfinitecond1}
    Assume that $\varphi,\psi\in \PSH(X,\theta)$ and
    \begin{equation}\label{eq:varphimassppsimass_greater}
        \int_X \theta_{\varphi}^n+\int_X \theta_{\psi}^n>\int_X \theta_{\varphi\lor\psi}^n.
    \end{equation}
    Then $\varphi\land \psi\in \PSH(X,\theta)$.
\end{proposition}
Thanks to the monotonicity theorem \cref{thm:mono}, we may also replace $\varphi\lor \psi$ on the right-hand side of \eqref{eq:varphimassppsimass_greater} by any $\gamma\in \PSH(X,\theta)$ such that $\varphi\lor \psi\preceq \gamma$.
\begin{proof}
    Without loss of generality, we may assume that $\varphi,\psi\leq 0$. For simplicity, we write
    \[
    \eta=P_{\theta}[\varphi\lor \psi].
    \]
    Note that due to our assumption \eqref{eq:varphimassppsimass_greater}, $\varphi,\psi\in \PSH(X,\theta)_{>0}$, and hence so is $\varphi\lor \psi$ thanks to the monotonicity theorem \cref{thm:mono}.
    By \cref{cor:Psendspotentialtomodel}, $\eta$ is a model potential, and $\int_X\theta_\eta^n=\int_X\theta_{\varphi\lor\psi}^n>0$.

    Take $C>0$ large enough, so that
    \begin{equation}\label{eq:intsumintbiggerinttemp1}
    \int_{\{\varphi>\eta-C\}} \theta_{\varphi}^n+\int_{\{\psi>\eta-C\}} \theta_{\psi}^n>\int_X \theta_{\eta}^n.
    \end{equation}
    This is possible thanks to \itemref{itm:prop:npp1:4}.
    Fix $C'>C$.
    Write
    \[
    \gamma_{C'}\coloneqq \Bigl(\varphi\lor (\eta-C')\Bigr)\land \Bigl(\psi\lor (\eta-C')\Bigr).
    \]
    Then observe that $\gamma_{C'}\sim \eta$, and
    \[
    \inf_{C'>C}\gamma_{C'}=\varphi\land \psi.
    \]
    Suppose, for contradiction, that $\varphi\land\psi\equiv -\infty$. Then
    \[
    \adjustlimits\lim_{C'\to \infty} \sup_X\gamma_{C'}=-\infty.
    \]
    Thanks to \cref{prop:model_sup_prop}, for each $C'>C$,
    \[
    \sup_{X} \gamma_{C'}=\sup_{\{\eta\neq -\infty\}} (\gamma_{C'}-\eta),
    \]
    since $\eta$ is a model potential.
    It follows that
    \begin{equation}\label{eq:limsupgammametatemp1}
        \adjustlimits \lim_{C'\to \infty} \sup_{\{\eta\neq -\infty\}}  (\gamma_{C'}-\eta)=-\infty.
    \end{equation}
    In particular, we can take $C'$ large enough so that
    \[
    \gamma_{C'}\leq \eta-C.
    \]
    For such $C'$, we have
    \[
    \begin{aligned}
        \int_X \theta_{\gamma_{C'}}^n=&\int_{\{\gamma_{C'}\leq \eta-C\}} \theta_{\gamma_{C'}}^n \\
        \leq & \int_{\left\{\varphi\lor (\eta-C')\leq \eta-C\right\}} \theta_{\varphi\lor (\eta-C')}^n +\int_{\left\{\psi\lor (\eta-C')\leq \eta-C \right\}} \theta_{\psi\lor (\eta-C')}^n\\
        =& 2\int_X \theta_{\eta}^n-\int_{\{\varphi>\eta-C\}}\theta_{\varphi}^n-\int_{\{\psi>\eta-C\}}\theta_{\psi}^n\\
        <& \int_X \theta_{\eta}^n,
    \end{aligned}
    \]
    where the second line follows from \cref{lma:rooftopMA}, the fourth line follows from \eqref{eq:intsumintbiggerinttemp1}. This contradicts the fact that $\gamma_{C'}\sim \eta$ in view of \cref{thm:mono}.
\end{proof}

\begin{proposition}\label{prop:landpresmodel}
    Let $\varphi,\psi\in \PSH(X,\theta)$. Assume that
    \begin{equation}\label{eq:varphipsimodel_landpsh}
        \varphi=P_{\theta}[\varphi],\quad \psi=P_{\theta}[\psi],\quad \varphi\land \psi\not\equiv -\infty.
    \end{equation}
    Then
    \begin{equation}\label{eq:Pthetaphilandpsi}
        P_{\theta}[\varphi\land \psi]=\varphi\land \psi.
    \end{equation}
\end{proposition}
\begin{proof}
    Observe that by \itemref{itm:prop:Pconc:3},
    \[
        P_{\theta}[\varphi\land \psi]\leq P_{\theta}[\varphi]=\varphi,\quad P_{\theta}[\varphi\land \psi]\leq P_{\theta}[\psi]=\psi.
    \]
    So the $\leq$ direction in \eqref{eq:Pthetaphilandpsi} holds. The reverse direction is trivial.
\end{proof}

\begin{lemma}\label{lma:landcorrectmass}
    Let $\varphi,\psi\in \PSH(X,\theta)$. Assume that \eqref{eq:varphipsimodel_landpsh} holds and
    \[
    \int_X \theta_{\varphi\land \psi}^n>0.
    \]
    Suppose that $\gamma,\eta\in \PSH(X,\theta)$ and satisfy
    \[
    \gamma\preceq \varphi,\quad \eta\preceq \psi,\quad \int_X \theta_{\gamma}^n=\int_X \theta_{\varphi}^n,\quad \int_X \theta_{\eta}^n=\int_X \theta_{\psi}^n.
    \]
    Then $\gamma\land \eta\not\equiv -\infty$ and
    \begin{equation}\label{eq:thetagammaeta_temp1}
    \int_X \theta_{\gamma\land \eta}^n=\int_X \theta_{\varphi\land \psi}^n.
    \end{equation}
\end{lemma}
Later on in \cref{chap:comp}, we shall reformulate and generalize this lemma in terms of the $P$-preorder, which is a more natural language to express the result. See \cref{prop:rooftopprePequiv}.
\begin{proof}
    Choose constants $A,B>0$ such that
    \[
        \gamma-A\leq \varphi,\quad \eta-B\leq \psi.
    \]
    Since
    \[
        (\gamma-A)\land(\eta-B)\sim \gamma\land\eta,
    \]
    the non-pluripolar masses of these two rooftop potentials are equal by the monotonicity theorem \cref{thm:mono}. Hence it suffices to prove the lemma after replacing $\gamma,\eta$ by $\gamma-A,\eta-B$. We can therefore assume that
    \[
        \gamma\leq \varphi,\quad \eta\leq \psi.
    \]

    \textbf{Step~1}. We first show that $\gamma\land \psi\not\equiv -\infty$ and
    \begin{equation}\label{eq:theta_gammalandpsi_mass_temp1}
    \int_X \theta_{\gamma\land \psi}^n=\int_X \theta_{\varphi\land \psi}^n.
    \end{equation}

    For any $a>1$, we define
    \[
    \gamma_a\coloneqq P_{\theta}\Bigl( a\gamma+(1-a)\varphi \Bigr).
    \]
    When $a=1$, we simply write $\gamma_1=\gamma$.

    Thanks to \cref{cor:pathoenvelopeeqmass}, for any $a\geq 1$, we have $\gamma_a\leq \varphi$ and
    \[
    \int_X \theta_{\gamma_a}^n=\int_X \theta_{\varphi}^n.
    \]
    By our assumption, for any $a\geq 1$, we have
    \[
    \int_X \theta_{\gamma_a}^n+\int_X \theta_{\varphi\land \psi}^n> \int_X \theta_{\varphi}^n,\quad \gamma_a\preceq \varphi,\quad \varphi\land \psi\preceq \varphi.
    \]
    So \cref{prop:landfinitecond1} gives
    \[
    \gamma_a\land \varphi\land \psi=\gamma_a\land \psi\in \PSH(X,\theta).
    \]
    In particular, for $a=1$, we find $\gamma\land \psi\not\equiv -\infty$. It remains to verify \eqref{eq:theta_gammalandpsi_mass_temp1}.

    For $a>1$, we have
    \[
    \gamma\geq a^{-1}\gamma_a+\left(1-a^{-1}\right)\varphi,
    \]
    thanks to the definition of $\gamma_a$ (cf. \eqref{eq:P_path_everywhere_ineq}). Hence
    \[
    \gamma\land \psi\geq a^{-1}\left( \gamma_a\land \psi\right)+\left(1-a^{-1}\right)\left(\varphi\land \psi\right).
    \]
    Therefore, by the monotonicity theorem \cref{thm:mono},
    \[
    \int_X \theta_{\gamma\land \psi}^n\geq \left(1-a^{-1}\right)^n \int_X \theta_{\varphi\land \psi}^n.
    \]
    Letting $a\to \infty$ and taking \cref{thm:mono} into account, \eqref{eq:theta_gammalandpsi_mass_temp1} follows.

    \textbf{Step~2}. We complete the proof.

    For any $a>1$, we let
    \[
    \eta_a\coloneqq P_{\theta}\Bigl( a\eta+(1-a)\psi \Bigr).
    \]
    When $a=1$, we set $\eta_1=\eta$.
    Thanks to \cref{cor:pathoenvelopeeqmass}, for each $a\geq 1$, we have $\eta_a\leq \psi$ and
    \[
    \int_X \theta_{\eta_a}^n=\int_X \theta_{\psi}^n.
    \]
    From Step~1, we have
    \[
    \int_X \theta_{\gamma\land\psi}^n+\int_X \theta_{\eta_a}^n=\int_X \theta_{\varphi\land \psi}^n+\int_X \theta_{\psi}^n>\int_X \theta_{\psi}^n.
    \]
    Applying \cref{prop:landfinitecond1} again, we find
    \[
    \gamma\land\psi\land \eta_a =\gamma\land \eta_a\in \PSH(X,\theta).
    \]
    When $a=1$, we find $\gamma\land \eta\in \PSH(X,\theta)$. It remains to prove \eqref{eq:thetagammaeta_temp1}.

    By definition of $\eta_a$, we have
    \[
    \eta\geq a^{-1}\eta_a+\left(1-a^{-1}\right)\psi.
    \]
    Therefore,
    \[
    \gamma\land \eta\geq a^{-1} (\gamma\land \eta_a)+\left(1-a^{-1}\right) \left( \gamma\land\psi \right).
    \]
    By the monotonicity theorem \cref{thm:mono},
    \[
    \int_X \theta_{\gamma\land \eta}^n\geq \left(1-a^{-1}\right)^n\int_X \theta_{\gamma\land \psi}^n.
    \]
    Letting $a\to \infty$ and applying \cref{thm:mono} again, we arrive at \eqref{eq:thetagammaeta_temp1}.
\end{proof}

There is an interesting diamond-like inequality regarding the non-pluripolar masses:
\begin{theorem}\label{thm:diamond}
    Assume that $\varphi,\psi\in \PSH(X,\theta)$ and $\varphi\land \psi\in \PSH(X,\theta)$. Then
    \begin{equation}\label{eq:diamond}
        \int_X \theta_{\varphi}^n+\int_X \theta_{\psi}^n\leq \int_X \theta_{\varphi\lor \psi}^n+\int_X \theta_{\varphi\land \psi}^n.
    \end{equation}
\end{theorem}
\begin{proof}
    We may assume that $\varphi,\psi\in \PSH(X,\theta)_{>0}$, as otherwise \eqref{eq:diamond} follows immediately from the monotonicity theorem \cref{thm:mono}.

    \textbf{Step~1}. We claim that it suffices to prove \eqref{eq:diamond} with $P_{\theta}[\varphi]$ and $P_{\theta}[\psi]$ in place of $\varphi$ and $\psi$.

    In fact, as $\varphi\land \psi\not \equiv -\infty$, we also have $P_{\theta}[\varphi]\land P_{\theta}[\psi]\not\equiv -\infty$. Moreover,
    \[
    \int_X \theta_{P_{\theta}[\varphi]}^n=\int_X \theta_{\varphi}^n,\quad \int_X \theta_{P_{\theta}[\psi]}^n=\int_X \theta_{\psi}^n
    \]
    by \cref{prop:Ppresmass}. On the other hand, as $C\to \infty$, the potentials
    \[
        \Bigl( (\varphi+C)\land V_{\theta} \Bigr) \lor \Bigl( (\psi+C)\land V_{\theta} \Bigr)
    \]
    converge to $P_{\theta}[\varphi]\lor P_{\theta}[\psi]$ almost everywhere. Therefore, thanks to \cref{cor:incseqnppcont} and \cref{thm:mono},
    \[
    \int_X \theta_{P_{\theta}[\varphi]\lor P_{\theta}[\psi]}^n=\int_X \theta_{\varphi\lor \psi}^n.
    \]
    Finally, we have
    \[
    \int_X \theta_{P_{\theta}[\varphi]\land P_{\theta}[\psi]}^n= \int_X \theta_{\varphi\land\psi}^n
    \]
    as well. When the left-hand side vanishes, this follows from \cref{thm:mono}. Otherwise, it follows from \cref{lma:landcorrectmass}.

    In particular, \eqref{eq:diamond} is equivalent to the corresponding result with $P_{\theta}[\varphi]$ and $P_{\theta}[\psi]$ in place of $\varphi$ and $\psi$.

    \textbf{Step~2}. We shall assume that $\varphi$ and $\psi$ are model potentials in the sequel.

    We write $\gamma=\varphi\lor \psi$ for simplicity. Note that $\gamma\leq 0$.

    For each $C>0$, we introduce
    \[
    \varphi_C=\varphi\lor (\gamma-C),\quad \psi_C=\psi\lor (\gamma-C).
    \]
    Note that $\varphi_C\sim \psi_C\sim \gamma$. Observe that \[
        \inf_{C>0} \varphi_C\land \psi_C=\varphi\land\psi.
    \]
    The $\geq$ direction is trivial, and the $\leq$ direction follows from the observation that the left-hand side is dominated by both $\varphi_C$ and $\psi_C$ for any $C>0$, and therefore by both $\varphi$ and $\psi$.

    For $C>0$, we compute
    \begin{equation}\label{eq:thetagammaminusvarphi_temp2}
    \begin{aligned}
    \theta_{\varphi_C}^n=& \mathds{1}_{\{\varphi>\gamma-C\}} \,\theta_{\varphi_C}^n+ \mathds{1}_{\{\varphi\leq \gamma-C\}} \,\theta_{\varphi_C}^n &\\
    =& \mathds{1}_{\{\varphi>\gamma-C\}} \,\theta_{\varphi}^n+ \mathds{1}_{\{\varphi\leq \gamma-C\}} \,\theta_{\varphi_C}^n & \textup{by \itemref{itm:prop:npp1:1}}\\
    =& \theta_{\varphi}^n+ \mathds{1}_{\{\varphi\leq \gamma-C\}} \,\theta_{\varphi_C}^n & \textup{by \cref{thm:Pvarphisupport}}.
    \end{aligned}
\end{equation}
    In the last line we use that $\varphi$ is model, so $\theta_\varphi^n$ is carried by $\{\varphi=0\}\subseteq\{\varphi>\gamma-C\}$.
    We also get a similar formula for $\theta_{\psi_C}^n$.

    Therefore, taking the monotonicity theorem \cref{thm:mono} into consideration, we find
    \begin{equation}\label{eq:intthetagammaminusvarphi_temp1}
    \int_X \theta_{\gamma}^n-\int_X \theta_{\varphi}^n=\int_{\{\varphi\leq \gamma-C\}} \theta_{\varphi_C}^n,\quad \int_X \theta_{\gamma}^n-\int_X \theta_{\psi}^n=\int_{\{\psi\leq \gamma-C\}} \theta_{\psi_C}^n.
    \end{equation}
    On the other hand, using \cref{lma:rooftopMA} and \eqref{eq:thetagammaminusvarphi_temp2}, we find
    \[
    \begin{aligned}
    \theta_{\varphi_C\land \psi_C}^n \leq & \mathds{1}_{\{\varphi_C\land \psi_C=\varphi_C\}} \,\theta_{\varphi_C}^n+\mathds{1}_{\{\varphi_C\land \psi_C=\psi_C\}} \,\theta_{\psi_C}^n \\
    \leq & \mathds{1}_{\{\varphi_C\land \psi_C=\varphi_C=0\}}\,\theta_{\varphi}^n+\mathds{1}_{\{\varphi_C\land \psi_C=\psi_C=0\}}\,\theta_{\psi}^n+\mathds{1}_{\{\varphi\leq \gamma-C\}}\,\theta_{\varphi_C}^n\\
    &+\mathds{1}_{\{\psi\leq \gamma-C\}}\,\theta_{\psi_C}^n.
    \end{aligned}
    \]
    In particular,
    \[
    \mathds{1}_{\{ \varphi_C\land \psi_C<0 \}}\, \theta_{\varphi_C\land \psi_C}^n \leq \mathds{1}_{\{\varphi\leq \gamma-C\}}\,\theta_{\varphi_C}^n+\mathds{1}_{\{\psi\leq \gamma-C\}}\,\theta_{\psi_C}^n.
    \]
    Integrating, we find
    \[
    \begin{aligned}
    \int_{\{ \varphi_C\land \psi_C=0 \}} \theta_{\varphi_C\land \psi_C}^n
    \geq & \int_X \theta_{\gamma}^n-\int_{\{\varphi\leq \gamma-C\}} \theta_{\varphi_C}^n-\int_{\{\psi\leq \gamma-C\}} \theta_{\psi_C}^n\\
    =& \int_X \theta_{\varphi}^n+\int_X \theta_{\psi}^n-\int_X \theta_{\gamma}^n,
    \end{aligned}
    \]
    where on the first line we used the fact that $\varphi_C\land \psi_C\sim \gamma$, while on the second line, we used \eqref{eq:intthetagammaminusvarphi_temp1}.

    Letting $C\to\infty$, and using \cref{cor:intcontlimsup}, we conclude \eqref{eq:diamond}.
\end{proof}


For later use, we also introduce a generalized envelope operator:
\begin{definition}\label{def:Pthetavarphif}
Given $\varphi\in\PSH(X,\theta)$ and a function $f\colon X\to[-\infty,\infty]$, we define
\[
P_\theta[\varphi](f)\coloneqq \uscsup \bigl\{\eta\in\PSH(X,\theta): \eta\leq f\textup{ quasi-everywhere},\ \eta\preceq\varphi\bigr\}.
\]
\index{$P_\theta[\varphi](f)$}
\end{definition}
In general, we cannot replace $\eta\leq f$ quasi-everywhere by everywhere.


\begin{example}\label{ex:envelopes:L625}
    When $f=0$, we have $P_\theta[\varphi](f)=P_{\theta}[\varphi]$.

    When $\varphi=V_{\theta}$, we have $P_\theta[V_{\theta}](f)=P_{\theta}(f)$.

    So \cref{def:Pthetavarphif} is a common generalization of \cref{def:Penv} and \cref{def:Pthetaf}.
\end{example}


\begin{proposition}\label{prop:DNTenvelope2}
    Let $\varphi\in\PSH(X,\theta)$ and $f\in C^{1,\overline{1}}(X)$ with $\varphi\leq f$. Then
    \[
    \theta^n_{P_\theta[\varphi](f)} = \mathds{1}_{\left\{P_\theta[\varphi](f)=f\right\}}\,\theta^n_f.
    \]
    In particular,
    \[
    \theta^n_{P_\theta[\varphi]} = \mathds{1}_{\left\{P_\theta[\varphi]=0\right\}}\,\theta^n.
    \]
\end{proposition}
This is a vast generalization of \cref{thm:Pvarphisupport}.
\begin{proof}
    By definition, we have
    \[
    P_\theta[\varphi](f)=P_\theta[\varphi]\Bigl(P_\theta(f)\Bigr)
    \]
    and $\varphi\leq P_\theta(f)$. By \cref{cor:Pvarphisupport} and \cref{cor:DNTenvelope1},
    \[
        \theta^n_{P_\theta[\varphi](f)} \leq \mathds{1}_{\left\{P_\theta[\varphi](f)=P_\theta(f)\right\}}\,\theta^n_{P_\theta(f)} = \mathds{1}_{\left\{P_\theta[\varphi](f)=f\right\}}\,\theta^n_f.
    \]
    The reverse inequality follows immediately from \cref{thm:DNTcontact}.
\end{proof}


\begin{proposition}\label{prop:DNTenvelope3}
    Let $\varphi\in\PSH(X,\theta)$ and $f\in C^{1,\overline{1}}(X)$ with $\varphi\leq f$. Then
    \[
    \begin{aligned}
        \mathds{1}_{\left\{\varphi=P_\theta(f)\right\}}\,\theta^n_{P_\theta(f)} =& \mathds{1}_{\left\{\varphi=P_\theta(f)\right\}}\,\theta^n_\varphi, \\
        \mathds{1}_{\left\{\varphi=P_\theta[\varphi](f)\right\}}\,\theta^n_{P_\theta[\varphi](f)} =& \mathds{1}_{\{\varphi=P_\theta[\varphi](f)\}}\,\theta^n_\varphi.
    \end{aligned}
    \]
\end{proposition}
\begin{proof}
    By \cref{cor:DNTenvelope1} and \cref{thm:DNTcontact}, we have
    \[
        \mathds{1}_{\{\varphi=P_\theta(f)\}}\,\theta^n_{P_\theta(f)}=\mathds{1}_{\{\varphi=P_\theta(f)=f\}}\,\theta^n_{f}=\mathds{1}_{\{\varphi=P_\theta(f)=f\}}\,\theta^n_{\varphi}\leq \mathds{1}_{\{\varphi=P_\theta(f)\}}\,\theta^n_{\varphi}.
    \]
    The reverse inequality follows from \cref{lma:lorMAineqe}. The first equality follows.

    As for the second, we compute using \cref{prop:DNTenvelope2} and \cref{thm:DNTcontact} that
    \[
    \begin{split}
    \mathds{1}_{\left\{\varphi=P_\theta[\varphi](f)\right\}}\,\theta^n_{P_\theta[\varphi](f)}=\mathds{1}_{\left\{\varphi=P_\theta[\varphi](f)=f\right\}}\,\theta^n_{f}\\
    =\mathds{1}_{\left\{\varphi=P_\theta[\varphi](f)=f\right\}}\,\theta_{\varphi}^n\leq \mathds{1}_{\{\varphi=P_\theta[\varphi](f)\}}\,\theta^n_\varphi.
    \end{split}
    \]
    Since $\varphi\leq P_{\theta}[\varphi](f)$, \cref{lma:lorMAineqe} gives the opposite inequality. The second equality follows.
\end{proof}

\subsection{Properties of the \texorpdfstring{$P$}{P}-envelope}
\label{subsec:Penv_prop}
Let $\theta,\theta_1,\theta_2$ be smooth closed real $(1,1)$-forms on $X$.

\begin{proposition}\label{prop:Penvbimero}
    Let $\pi\colon Y\rightarrow X$ be a proper bimeromorphic morphism from a Kähler manifold $Y$ to $X$. Then for any $\varphi\in \PSH(X,\theta)$, we have
    \[
        P_{\pi^*\theta}[\pi^*\varphi]=\pi^*P_{\theta}[\varphi].
    \]
    In particular, a potential $\varphi\in \PSH(X,\theta)_{>0}$ is model if and only if $\pi^*\varphi\in \PSH(Y,\pi^*\theta)_{>0}$ is model.
\end{proposition}
\begin{proof}
    By \cref{prop:PSHpullbij}, every $\chi\in\PSH(Y,\pi^*\theta)$ is of the form $\chi=\pi^*\psi$ for a unique $\psi\in\PSH(X,\theta)$. Since $\pi$ is surjective, $\chi\leq 0$ if and only if $\psi\leq 0$. Moreover,
    \[
        \chi\preceq \pi^*\varphi \Longleftrightarrow \psi\preceq \varphi.
    \]
    Thus the candidates defining $P_{\pi^*\theta}[\pi^*\varphi]$ are exactly the pull-backs of the candidates defining $P_\theta[\varphi]$, and hence
    \[
        P_{\pi^*\theta}[\pi^*\varphi]=\pi^*P_\theta[\varphi].
    \]
    Finally, the non-pluripolar mass is preserved under bimeromorphic pull-back by \cref{prop:nppmassinv}. Therefore $\varphi$ is model of positive mass if and only if $\pi^*\varphi$ is.
\end{proof}

We have the following concavity property of the $P$-envelope.
\begin{proposition}\label{prop:Pconc}
    \leavevmode
    \begin{enumerate}
        \item\label{itm:prop:Pconc:1} Suppose that $\varphi\in \PSH(X,\theta)$ and $\lambda\in \mathbb{R}_{>0}$. Then
            \[
                P_{\lambda\theta}[\lambda\varphi]=\lambda P_{\theta}[\varphi].
            \]
        \item\label{itm:prop:Pconc:2} Suppose that $\varphi_1\in \PSH(X,\theta_1)$ and $\varphi_2\in \PSH(X,\theta_2)$. Then
            \[
                P_{\theta_1+\theta_2}[\varphi_1+\varphi_2]\geq P_{\theta_1}[\varphi_1]+P_{\theta_2}[\varphi_2].
            \]
        \item\label{itm:prop:Pconc:3} Suppose that $\varphi,\psi\in \PSH(X,\theta)$ and $\varphi\preceq \psi$. Then
            \[
                P_{\theta}[\varphi]\leq P_{\theta}[\psi].
            \]
    \end{enumerate}
\end{proposition}
\begin{proof}
    \ref{itm:prop:Pconc:1} This follows by rescaling the candidates in the definition of the $P$-envelope.

    \ref{itm:prop:Pconc:2} Suppose that $\psi_1\in \PSH(X,\theta_1)$ and $\psi_2\in \PSH(X,\theta_2)$ satisfy
    \[
        \psi_i\leq 0,\quad \psi_i\preceq \varphi_i
    \]
    for $i=1,2$. Then
    \[
        \psi_1+\psi_2\leq 0,\quad \psi_1+\psi_2\preceq \varphi_1+\varphi_2.
    \]
    It follows from \eqref{eq:Pthetavarphi} that
    \[
        \psi_1+\psi_2\leq P_{\theta_1+\theta_2}[\varphi_1+\varphi_2].
    \]
    Since $\psi_1$ and $\psi_2$ are arbitrary, we conclude.

    \ref{itm:prop:Pconc:3} This follows by inclusion of the defining candidate sets.
\end{proof}











\begin{proposition}\label{prop:decseqmodel}
    Let $(\varphi_j)_{j\in I}$ be a decreasing net of potentials in $\PSH(X,\theta)$ satisfying $P_{\theta}[\varphi_j]=\varphi_j$ for each $j\in I$. Set $\varphi\coloneqq \inf_{j\in I} \varphi_j$. Then $P_{\theta}[\varphi]=\varphi$.
\end{proposition}
\begin{proof}
    Since $P_{\theta}[\varphi_j]=\varphi_j$, we have $\sup_X\varphi_j=0$. By \cref{cor:decreasingnetL1conv}, we have $\varphi\in \PSH(X,\theta)$, and $\varphi\leq 0$. Therefore, for each $j\in I$,
    \[
        \varphi\leq P_{\theta}[\varphi]\leq P_{\theta}[\varphi_j] =\varphi_j.
    \]
    Therefore, $\varphi=P_{\theta}[\varphi]$.
\end{proof}


\begin{proposition}\label{prop:vol_limit_model}
    Let $(\epsilon_i)_{i\in I}$ be a decreasing net in $\mathbb{R}_{\geq 0}$ with limit $0$. Take a Kähler form $\omega$ on $X$.
    Consider a decreasing net $(\varphi_i)_{i\in I}$ with $\varphi_i\in \PSH(X,\theta+\epsilon_i \omega)$ satisfying
    \begin{equation}\label{eq:Palmostmodeltemp}
        P_{\theta+\epsilon_i\omega}\left[\varphi_i\right]=\varphi_i
    \end{equation}
    for every $i\in I$. Set $\varphi\coloneqq\inf_{i\in I}\varphi_i$. Then $P_{\theta}[\varphi]=\varphi$, and
    \begin{equation}\label{eq:massmodeldec}
    \lim_{i\in I}\int_X \left(\theta + \epsilon_i \omega\right)^n_{\varphi_i}=\int_X \theta_{\varphi}^n.
    \end{equation}
    Moreover, if $\int_X \theta_{\varphi}^n>0$, then for any prime divisor $E$ over $X$, we have
    \begin{equation}\label{eq:Lelongcontdecseq}
    \lim_{i\in I} \nu\left(\varphi_i,E\right)=\nu(\varphi,E).
    \end{equation}
\end{proposition}
Note that both \eqref{eq:massmodeldec} and \eqref{eq:Lelongcontdecseq} fail without the assumption \eqref{eq:Palmostmodeltemp}.

\begin{proof}
        We first observe that $\varphi\in \PSH(X,\theta)$. In fact, for each $i_0\in I$, the tail $(\varphi_i)_{i\geq i_0}$ lies in $\PSH(X,\theta+\epsilon_{i_0}\omega)$ and decreases to $\varphi$. Since $\sup_X\varphi_i=0$ by \eqref{eq:Palmostmodeltemp}, the limit is not identically $-\infty$. Thus $\varphi\in \PSH(X,\theta+\epsilon_{i_0}\omega)$ for every $i_0\in I$. Since $\epsilon_{i_0}\to 0$, it follows that $\varphi\in \PSH(X,\theta)$.
        For every $i\in I$ we have
        \[
        P_{\theta}[\varphi]\leq P_{\theta+\epsilon_i\omega}[\varphi]\leq P_{\theta+\epsilon_i\omega}[\varphi_i]=\varphi_i.
        \]
        Therefore, $P_{\theta}[\varphi]\leq \varphi$. The reverse inequality is trivial. Hence $P_{\theta}[\varphi]=\varphi$.

        By the monotonicity theorem \cref{thm:mono} and the multi-linearity of the non-pluripolar product \itemref{itm:prop:npp1:6}, we have
        \[
            \varliminf_{i\in I} \int_X \left(\theta + \epsilon_i \omega\right)^n_{\varphi_i} \geq  \varliminf_{i\in I} \int_X \left(\theta + \epsilon_i \omega\right)^n_{\varphi}  = \int_X \theta_{\varphi}^n.
        \]
        We now argue the reverse inequality.

        Fix $i_0\in I$. Since passing to the tail $i\geq i_0$ does not change the limsup, we have
        \[
            \begin{aligned}
                \varlimsup_{i\in I} \int_X \left(\theta + \epsilon_i \omega\right)^n_{\varphi_i}=&\varlimsup_{i\geq i_0} \int_{\left\{\varphi_i =0\right\}} \left(\theta + \epsilon_i \omega\right)^n_{\varphi_i} \\
                \leq & \varlimsup_{i\geq i_0} \int_{\left\{\varphi_i =0\right\}} \left(\theta + \epsilon_{i_0} \omega\right)^n_{\varphi_i}\\
                \leq & \int_{\left\{\varphi =0\right\}} \left(\theta + \epsilon_{i_0} \omega\right)^n_{\varphi},
            \end{aligned}
        \]
    where in the first line we used \eqref{eq:Palmostmodeltemp} and \cref{thm:Pvarphisupport}, and in the last line we used the fact that $\varphi_i \searrow \varphi$ on the tail $i\geq i_0$, and on this tail, $\varphi_i \to \varphi$ in $L^1$ as well by \cref{cor:decreasingnetL1conv}, and therefore, \cref{cor:intcontlimsup} is applicable. Taking the limit with respect to $i_0$, we arrive at the desired conclusion:
    \[
    \varlimsup_{i\in I} \int_{X} \left(\theta + \epsilon_i \omega \right)^n_{\varphi_i} \leq \varliminf_{i_0\in I}\int_{\{\varphi =0\}} \left(\theta + \epsilon_{i_0} \omega \right)^n_{\varphi}  =  \int_{\{\varphi =0\}} \theta^n_{\varphi} \leq \int_{X} \theta^n_{\varphi},
    \]
    where the equality follows from the multi-linearity of the non-pluripolar product \itemref{itm:prop:npp1:6} and the last inequality follows from \cref{thm:Pvarphisupport}.
    This finishes the proof of \eqref{eq:massmodeldec}.

    It remains to argue \eqref{eq:Lelongcontdecseq}.
    Fix $i_0\in I$. By \cref{prop:Lelongnumberupperbound}, there is a constant $B>0$ such that
    \[
    \nu(\psi,E)\leq B,\quad \psi\in \PSH(X,\theta+\epsilon_{i_0}\omega).
    \]
    By \cref{lma:pathoenvelope} and \eqref{eq:massmodeldec}, for any $\delta \in (0,1)$ there exists $i_{\delta}\geq i_0$ such that for all $i\geq i_{\delta}$ there exists $\psi_i \in \PSH(X,\theta+\epsilon_i\omega)$ satisfying $(1-\delta) \varphi_i + \delta \psi_i \leq \varphi$. This implies that for $i\geq i_{\delta}$ we have
    \[
    (1-\delta)\nu(\varphi_i,E) + \delta\nu(\psi_i,E) \geq \nu(\varphi,E) \geq \nu(\varphi_i,E).
    \]
    Since $\psi_i\in \PSH(X,\theta+\epsilon_{i_0}\omega)$ for $i\geq i_{\delta}$, we get
    \[
        \varliminf_{i\in I}\nu(\varphi_i,E)\geq \frac{\nu(\varphi,E)-\delta B}{1-\delta}.
    \]
    Letting $\delta\searrow 0$ gives $\varliminf_{i\in I}\nu(\varphi_i,E)\geq \nu(\varphi,E)$. The reverse inequality follows from $\varphi\leq \varphi_i$, and hence we conclude \eqref{eq:Lelongcontdecseq}.
\end{proof}








\begin{corollary}\label{cor:varphiperturbtheta}
    Let $\varphi\in \PSH(X,\theta)_{>0}$ be a model potential. Let $\omega$ be a Kähler form on $X$.
    Then
    \[
    \varphi=\inf_{\epsilon>0} P_{\theta+\epsilon\omega}[\varphi].
    \]
\end{corollary}
\begin{proof}
    Set
    \[
        \eta\coloneqq\inf_{\epsilon>0}P_{\theta+\epsilon\omega}[\varphi].
    \]
    Since $\varphi$ is model, $\varphi=P_{\theta}[\varphi]\leq P_{\theta+\epsilon\omega}[\varphi]$ for every $\epsilon>0$, hence $\varphi\leq\eta\leq 0$.

    Next we prove the reverse direction.
    Applying \cref{prop:vol_limit_model} to the decreasing net $(P_{\theta+\epsilon\omega}[\varphi])_{\epsilon>0}$ and using \cref{prop:Ppresmass}, we get
    \[
        P_\theta[\eta]=\eta,\quad \int_X\theta_\eta^n=\int_X\theta_\varphi^n.
    \]
    Thus $\eta$ is a candidate in the characterization \cref{thm:Pvarphidiffdef} of $P_\theta[\varphi]$. Hence $\eta\leq P_\theta[\varphi]=\varphi$.
\end{proof}

\begin{proposition}\label{prop:incnetmodel}
    Let $(\varphi_i)_{i\in I}$ be an increasing net of potentials in $\PSH(X,\theta)_{>0}$ uniformly bounded from above. Let $\varphi\coloneqq \uscsup[i\in I]\varphi_i$. Then
    \[
    \uscsup[i\in I] P_{\theta}[\varphi_i]=P_{\theta}[\varphi].
    \]

    In particular, if $\varphi_i$ is model for all $i\in I$, then so is $\varphi$.
\end{proposition}
\begin{proof}
We may assume that $I$ is infinite since otherwise, there is nothing to prove.

    We write
    \[
    \eta\coloneqq \uscsup[i\in I] P_{\theta}[\varphi_i].
    \]
    Then it is clear that $\eta\leq P_{\theta}[\varphi]$.

    By \cref{cor:incseqnppcont}, we have
    \[
        \lim_{i\in I} \int_X\theta_{\varphi_i}^n=\int_X\theta_{\varphi}^n>0.
    \]
    We may assume that for each $i\in I$, $\int_X\theta_{\varphi_i}^n<\int_X\theta_{\varphi}^n$, since otherwise $P_{\theta}[\varphi_i]=P_{\theta}[\varphi]$ by \cref{thm:Pvarphidiffdef}, and the same holds for all $j\geq i$; in that case there is nothing to prove.

    So by \cref{lma:pathoenvelope}, we can find a decreasing net $\epsilon_i\searrow 0$ ($i\in I$) with $\epsilon_i\in (0,1)$ and $\psi_i\in \PSH(X,\theta)$ ($i\in I$) such that for all $i\in I$,
    \[
        (1-\epsilon_i) \varphi+\epsilon_i\psi_i\leq \varphi_i.
    \]
    We may, for example, take
    \begin{equation}\label{eq:epsilon_inetchoice}
        \epsilon_i^{-1}=\frac{1}{2}\left( 1+\left(\frac{\int_X \theta_{\varphi}^n}{\int_X \theta_{\varphi}^n-\int_X \theta_{\varphi_i}^n}\right)^{1/n} \right).
    \end{equation}


    By the concavity property of the envelope \cref{prop:Pconc}, we have
    \[
        P_{\theta}[\varphi]+\epsilon_i P_{\theta}[\psi_i]\leq (1-\epsilon_i) P_{\theta}[\varphi]+\epsilon_i P_{\theta}[\psi_i]\leq P_{\theta}[\varphi_i]\leq\eta.
    \]
    Each $P_{\theta}[\psi_i]$ has supremum $0$, hence the family is bounded in $L^1$ by \cref{prop:L1compa}. Thus $\epsilon_iP_{\theta}[\psi_i]\to 0$ in $L^1$. The previous inequality then gives $P_{\theta}[\varphi]\leq \eta$ almost everywhere, and hence everywhere by \cref{prop:pshfuncdetdense}.
\end{proof}




\subsection{Relative full mass classes}\label{subsec:fullmass}
Let $\theta$ be a smooth closed real $(1,1)$-form on $X$ representing a big cohomology class. Fix a model potential $\phi\in \PSH(X,\theta)_{>0}$.



\begin{definition}\label{def:envelopes:L906}
    We define
    \[
        \begin{aligned}
        \PSH(X,\theta;\phi)\coloneqq &\left\{\eta\in \PSH(X,\theta): \eta\preceq \phi\right\},\\
        \mathcal{E}^{\infty}(X,\theta;\phi)\coloneqq &\left\{\eta\in \PSH(X,\theta): \eta\sim \phi\right\},\\
        \mathcal{E}(X,\theta;\phi)\coloneqq &\left\{\eta\in \PSH(X,\theta;\phi): \int_X \theta_{\eta}^n=\int_X \theta_{\phi}^n\right\},\\
        \mathcal{E}^1(X,\theta;\phi)\coloneqq &\left\{\eta\in \mathcal{E}(X,\theta;\phi): \int_X |\phi-\eta|\,\theta_{\eta}^n<\infty\right\}.
        \end{aligned}
    \]
    \index{$\PSH(X,\theta;\phi)$}\index{$\mathcal{E}^{\infty}(X,\theta;\phi)$}\index{$\mathcal{E}(X,\theta;\phi)$} \index{$\mathcal{E}^1(X,\theta;\phi)$}
    Potentials in the last three classes are said to have \emph{relatively minimal singularities}\index{potential!potential with relatively minimal singularities}, \emph{full mass}\index{potential!potential with relative full mass} and \emph{finite energy}\index{potential!potential with relative finite energy} relative to $\phi$ respectively.
\end{definition}
We have the following inclusions:
\begin{equation}\label{eq:energyclassinc}
\mathcal{E}^{\infty}(X,\theta;\phi)\subseteq \mathcal{E}^1(X,\theta;\phi)\subseteq \mathcal{E}(X,\theta;\phi)\subseteq \PSH(X,\theta;\phi).
\end{equation}
The only non-trivial part is the first inclusion, which follows from the monotonicity theorem \cref{thm:mono}.

\begin{remark}\label{rmk:intwelldef}
    Note that this integral
    \[
    \int_X |\phi-\eta|\,\theta_{\eta}^n
    \]
    is defined: The locus where $\phi-\eta$ is undefined is a pluripolar set, while the product $\theta_{\eta}^n$ puts no mass on pluripolar sets (\cref{prop:npp1}).

    Similar remarks apply when we talk about similar integrals in the sequel.
\end{remark}


When $\phi=V_{\theta}$, we usually write
\[
    \begin{aligned}
        \mathcal{E}^{\infty}(X,\theta;V_{\theta})=& \mathcal{E}^{\infty}(X,\theta),\\
        \mathcal{E}(X,\theta;V_{\theta})=& \mathcal{E}(X,\theta),\\
        \mathcal{E}^{1}(X,\theta;V_{\theta})=& \mathcal{E}^{1}(X,\theta).
    \end{aligned}
\]
Potentials in the three classes are said to have \emph{minimal singularities}\index{potential!potential with minimal singularities}, \emph{full mass}\index{potential!potential with full mass} and \emph{finite energy}\index{potential!potential with finite energy} respectively. The relation \eqref{eq:energyclassinc} can be written as
\[
\mathcal{E}^{\infty}(X,\theta)\subseteq \mathcal{E}^{1}(X,\theta) \subseteq \mathcal{E}(X,\theta)
\]
in this case.

The $P$-envelope can be used to characterize the full mass classes:
\begin{proposition}\label{prop:fullmassP}
    Let $\varphi\in \PSH(X,\theta)$. Then the following are equivalent:
    \begin{enumerate}
        \item\label{itm:prop:fullmassP:1} $\varphi\in \mathcal{E}(X,\theta;\phi)$;
        \item\label{itm:prop:fullmassP:2} $P_{\theta}[\varphi]=\phi$.
    \end{enumerate}
\end{proposition}
\begin{proof}
    \ref{itm:prop:fullmassP:2} $\implies$ \ref{itm:prop:fullmassP:1}. Assume \ref{itm:prop:fullmassP:2}. By \cref{prop:Ppresmass},
    \[
    \int_X\theta_\varphi^n=\int_X\theta_{P_\theta[\varphi]}^n=\int_X\theta_\phi^n.
    \]
    Moreover, since the mass is positive, \cref{thm:Pvarphidiffdef} implies that
    \[
        \varphi\preceq P_\theta[\varphi]=\phi.
    \]
    Thus $\varphi\in\mathcal E(X,\theta;\phi)$.

    \ref{itm:prop:fullmassP:1} $\implies$ \ref{itm:prop:fullmassP:2}. Assume \ref{itm:prop:fullmassP:1}. Then $\varphi\in \PSH(X,\theta)_{>0}$.
    Note that $\phi$ is a candidate of $P_{\theta}[\varphi]$ as in \eqref{eq:Penvdef}. So $P_{\theta}[\varphi]\geq \phi$. The reverse inequality is trivial.
\end{proof}

We have the following comparison principle.
\begin{proposition}[Comparison principle]\label{prop:compa_princ}\index{theorem!comparison principle}
    Fix $j\in \{0,\ldots,n\}$.
    Let $\theta_1,\ldots,\theta_j$ be closed smooth real $(1,1)$-forms on $X$ and $\psi_i\in \PSH(X,\theta_i)$ for $i=1,\ldots,j$. Suppose that $\varphi,\psi\in \mathcal{E}(X,\theta;\phi)$. Then
    \begin{equation}\label{eq:comp_princi1}
    \int_{\{\varphi<\psi\}} \theta_{\varphi}^{n-j} \wedge \theta_{1,\psi_1}\wedge \cdots \wedge \theta_{j,\psi_j}\geq \int_{\{\varphi<\psi\}} \theta_{\psi}^{n-j} \wedge \theta_{1,\psi_1}\wedge \cdots \wedge \theta_{j,\psi_j}.
    \end{equation}
    As a consequence, we also have
    \begin{equation}\label{eq:comp_princi12}
    \int_{\{\varphi\leq \psi\}} \theta_{\varphi}^{n-j} \wedge \theta_{1,\psi_1}\wedge \cdots \wedge \theta_{j,\psi_j}\geq \int_{\{\varphi\leq\psi\}} \theta_{\psi}^{n-j} \wedge \theta_{1,\psi_1}\wedge \cdots \wedge \theta_{j,\psi_j}.
    \end{equation}
\end{proposition}
\begin{proof}
    Observe that $\varphi\lor \psi\in \mathcal{E}(X,\theta;\phi)$, as a consequence of the monotonicity theorem \cref{thm:mono}.
    Using the plurifine locality of the non-pluripolar product \itemref{itm:prop:npp1:1}, we compute
    \[
    \begin{aligned}
    &\int_X \theta_{\varphi\lor \psi}^{n-j}\wedge \theta_{1,\psi_1}\wedge \cdots \wedge \theta_{j,\psi_j}\\
    \geq & \int_{\{\varphi>\psi\} } \theta_{\varphi}^{n-j}\wedge \theta_{1,\psi_1}\wedge \cdots \wedge \theta_{j,\psi_j}+\int_{\{\varphi<\psi\} } \theta_{\psi}^{n-j}\wedge \theta_{1,\psi_1}\wedge \cdots \wedge \theta_{j,\psi_j}\\
    =& \int_X  \theta_{\varphi}^{n-j}\wedge \theta_{1,\psi_1}\wedge \cdots \wedge \theta_{j,\psi_j}-\int_{\{\varphi\leq \psi\} } \theta_{\varphi}^{n-j}\wedge \theta_{1,\psi_1}\wedge \cdots \wedge \theta_{j,\psi_j}\\
    &+\int_{\{\varphi<\psi\} } \theta_{\psi}^{n-j}\wedge \theta_{1,\psi_1}\wedge \cdots \wedge \theta_{j,\psi_j}\\
    =&\int_X  \theta_{\varphi\lor \psi}^{n-j}\wedge \theta_{1,\psi_1}\wedge \cdots \wedge \theta_{j,\psi_j}-\int_{\{\varphi\leq \psi\} } \theta_{\varphi}^{n-j}\wedge \theta_{1,\psi_1}\wedge \cdots \wedge \theta_{j,\psi_j}\\
    &+\int_{\{\varphi<\psi\} } \theta_{\psi}^{n-j}\wedge \theta_{1,\psi_1}\wedge \cdots \wedge \theta_{j,\psi_j},
    \end{aligned}
    \]
    where in the last step, we applied \cref{prop:Ppresmass} to $\varphi$ and $\varphi\lor \psi$. Therefore,
    \[
    \begin{aligned}
    \int_{\{\varphi\leq \psi\}}& \theta_{\varphi}^{n-j} \wedge \theta_{1,\psi_1}\wedge \cdots \wedge \theta_{j,\psi_j}\\
    &\geq \int_{\{\varphi<\psi\}} \theta_{\psi}^{n-j} \wedge \theta_{1,\psi_1}\wedge \cdots \wedge \theta_{j,\psi_j}.
    \end{aligned}
    \]
    Applying the preceding inequality to $\varphi+\epsilon$ in place of $\varphi$, and using $\theta_{\varphi+\epsilon}=\theta_\varphi$, we get
    \[
        \int_{\{\varphi+\epsilon\leq\psi\}}\theta_\varphi^{n-j}\wedge \theta_{1,\psi_1}\wedge\cdots\wedge\theta_{j,\psi_j} \geq \int_{\{\varphi+\epsilon<\psi\}}\theta_\psi^{n-j}\wedge \theta_{1,\psi_1}\wedge\cdots\wedge\theta_{j,\psi_j}.
    \]
    As $\epsilon\searrow 0$, by monotone convergence theorem, we obtain
    \[
    \int_{\{\varphi<\psi\}}\theta_\varphi^{n-j}\wedge \theta_{1,\psi_1}\wedge\cdots\wedge\theta_{j,\psi_j} \geq \int_{\{\varphi<\psi\}}\theta_\psi^{n-j}\wedge \theta_{1,\psi_1}\wedge\cdots\wedge\theta_{j,\psi_j},
    \]
    which proves \eqref{eq:comp_princi1}.

    Finally, the \eqref{eq:comp_princi12} follows by subtracting \eqref{eq:comp_princi1} (with the role of $\varphi$ and $\psi$ switched) from the total mixed mass equality
    \[
        \int_X\theta_\varphi^{n-j}\wedge \theta_{1,\psi_1}\wedge\cdots\wedge\theta_{j,\psi_j}= \int_X\theta_\psi^{n-j}\wedge \theta_{1,\psi_1}\wedge\cdots\wedge\theta_{j,\psi_j},
    \]
    which itself follows from \cref{prop:Ppresmass} and \cref{prop:fullmassP}.
\end{proof}

The full mass potentials are essential in resolving the Monge--Ampère equations. We recall the following two Aubin--Yau type theorems. Recall that a measure $\mu$ on $X$ is \emph{non-pluripolar}\index{non-pluripolar measure} if it is a Radon measure and $\mu(K)=0$ for each pluripolar set $K\subseteq X$.\footnote{Here we identify as usual a Radon measure with its completion.}
\begin{theorem}\label{thm:Yau_rel}
    Let $\mu$ be a non-pluripolar measure on $X$ with $\mu(X)=\int_X \theta_{\phi}^n$. Then there is a unique $\varphi\in \mathcal{E}(X,\theta;\phi)$ such that
    \[
    \theta_{\varphi}^n=\mu,\quad \sup_X \varphi=0.
    \]
\end{theorem}


\begin{theorem}\label{thm:AY_rel}
    Fix $\lambda>0$. Let $\mu$ be a non-pluripolar measure on $X$ with $\mu(X)>0$. Then there is a unique $\varphi\in \mathcal{E}(X,\theta;\phi)$ such that
    \[
    \theta_{\varphi}^n=\mathrm{e}^{\lambda\varphi}\mu.
    \]
    Furthermore, for any $\psi\in \mathcal{E}(X,\theta;\phi)$ satisfying
    \[
    \theta_{\psi}^n\geq \mathrm{e}^{\lambda\psi}\mu,
    \]
    we have $\varphi\geq \psi$.
\end{theorem}
The proofs are beyond the scope of this book; we refer to \cite{DDNL19log}.

In order to handle the finite energy classes, it is convenient to introduce the following quantity:
\begin{definition}\label{def:MAenergy}
    We define the \emph{Monge--Amp\`ere energy}\index{Monge--Amp\`ere energy} $E^{\phi}_{\theta}\colon \mathcal{E}^{\infty}(X,\theta;\phi)\rightarrow \mathbb{R}$\index{$E^{\phi}_{\theta}$} as follows
    \begin{equation}\label{eq:Edefbdd}
        E^{\phi}_{\theta}(\varphi)\coloneqq \frac{1}{n+1}\sum_{j=0}^n \int_X (\varphi-\phi)\, \theta_{\varphi}^j\wedge \theta_{\phi}^{n-j}.
    \end{equation}


    More generally, we extend $E_{\theta}^{\phi}$ to a functional $E^{\phi}_{\theta}\colon \PSH(X,\theta;\phi)\rightarrow [-\infty,\infty)$ as follows
    \begin{equation}\label{eq:Eextendgeneral}
        E^{\phi}_{\theta}(\varphi)\coloneqq \inf\left\{E^{\phi}_{\theta}(\psi):\psi\in \mathcal{E}^{\infty}(X,\theta;\phi),\varphi\leq \psi\right\}.
    \end{equation}

     We write $E_{\theta}$\index{$E_{\theta}$} instead of $E_{\theta}^{\phi}$ when $\phi=V_{\theta}$.
\end{definition}
The set in \eqref{eq:Eextendgeneral} is non-empty, since $\phi+\sup_X \varphi$ is a candidate. Therefore, $E^{\phi}_{\theta}(\varphi)<\infty$.


Note that
\begin{equation}\label{eq:Ephiaddcst}
E^{\phi}_{\theta}(\varphi+C)=E^{\phi}_{\theta}(\varphi)+C\int_X \theta_{\phi}^n
\end{equation}
for any $\varphi\in \PSH(X,\theta;\phi)$ and $C\in \mathbb{R}$.



\begin{lemma}\label{lma:EmonoEinfty}
    The functional $E_{\theta}^{\phi}\colon \mathcal{E}^{\infty}(X,\theta;\phi)\rightarrow \mathbb{R}$ is increasing. In particular, the extended definition \eqref{eq:Eextendgeneral} agrees with \eqref{eq:Edefbdd} on $\mathcal{E}^{\infty}(X,\theta;\phi)$, and is increasing as well.
\end{lemma}
\begin{proof}
    Let $\varphi,\psi\in \mathcal{E}^{\infty}(X,\theta;\phi)$, we have
    \[
    \begin{aligned}
    &(n+1)E^{\phi}_{\theta}(\varphi)-(n+1)E^{\phi}_{\theta}(\psi)-\sum_{j=0}^n \int_X (\varphi-\psi)\,\theta_{\varphi}^j\wedge \theta_{\psi}^{n-j} \\
    =& \sum_{j=0}^n \int_X (\varphi-\phi)\, \theta_{\varphi}^j\wedge \theta_{\phi}^{n-j}-\sum_{j=0}^n \int_X (\psi-\phi)\, \theta_{\psi}^j\wedge \theta_{\phi}^{n-j}-\sum_{j=0}^n \int_X (\varphi-\psi)\,\theta_{\varphi}^j\wedge \theta_{\psi}^{n-j}\\
    =& \sum_{j=0}^n \int_X (\varphi-\psi)\,\left( \theta_{\varphi}^j\wedge \theta_{\phi}^{n-j}-\theta_{\varphi}^j\wedge \theta_{\psi}^{n-j} \right)\\
    &+\sum_{j=0}^n \int_X (\psi-\phi)\,\left( \theta_{\varphi}^j\wedge \theta_{\phi}^{n-j}-\theta_{\psi}^j\wedge \theta_{\phi}^{n-j} \right)\\
    =& \sum_{\substack{j+a+b=n-1,\\ j,a,b\geq 0}} (n-j)\int_X (\psi-\phi)\,\left( \theta_{\phi}^a\wedge \theta_{\psi}^{b+1}\wedge \theta_{\varphi}^j- \theta_{\phi}^a\wedge \theta_{\psi}^{b}\wedge \theta_{\varphi}^{j+1}\right)\\
    &+\sum_{\substack{j+a+b=n-1,\\ j,a,b\geq 0}} (n-j)\int_X (\psi-\phi)\,\left( \theta_{\phi}^j\wedge \theta_{\varphi}^{a+1}\wedge \theta_{\psi}^b- \theta_{\phi}^j\wedge \theta_{\varphi}^{a}\wedge \theta_{\psi}^{b+1}\right)\\
    =&0,
    \end{aligned}
    \]
    where the third equality follows from the integration by parts formula. See \cite{Xia19b, Lu21} for the proof in the context of non-pluripolar products.

    In other words,
    \begin{equation}\label{eq:int_by_part_E_temp1}
    E^{\phi}_{\theta}(\varphi)-E^{\phi}_{\theta}(\psi)=\frac{1}{n+1}\sum_{j=0}^n \int_X(\varphi-\psi)\,\theta_{\varphi}^j\wedge \theta_{\psi}^{n-j}.
    \end{equation}
    The monotonicity follows.
\end{proof}


\begin{lemma}\label{lma:comp_int_ene}
    Let $\varphi,\psi\in \mathcal{E}^{\infty}(X,\theta;\phi)$. Then
    \begin{equation}\label{eq:comp_int_ene}
    \int_X (\varphi-\psi)\,\theta_{\psi}^n \geq E_{\theta}^{\phi}(\varphi)-E_{\theta}^{\phi}(\psi)\geq \int_X (\varphi-\psi)\,\theta_{\varphi}^n.
    \end{equation}
\end{lemma}
\begin{proof}
    Thanks to \eqref{eq:Ephiaddcst}, we may assume that $\varphi\geq \psi$.

       As we have seen in \eqref{eq:int_by_part_E_temp1}, the middle term can be written as
       \[
       \frac{1}{n+1}\sum_{j=0}^n \int_X(\varphi-\psi)\,\theta_{\varphi}^j\wedge \theta_{\psi}^{n-j}.
       \]
       We claim that for each individual $j$, the inequality \eqref{eq:comp_int_ene} holds:
       \[
       \int_X (\varphi-\psi)\,\theta_{\psi}^n \geq \int_X(\varphi-\psi)\,\theta_{\varphi}^j\wedge \theta_{\psi}^{n-j} \geq \int_X (\varphi-\psi)\,\theta_{\varphi}^n.
       \]
       Thanks to \cref{prop:compa_princ}, we have
       \[
       \begin{aligned}
       \int_X(\varphi-\psi)\,\theta_{\varphi}^j\wedge \theta_{\psi}^{n-j}=&\int_0^{\infty} \int_{\{\varphi>\psi+t\}} \theta_{\varphi}^j\wedge \theta_{\psi}^{n-j}\,\mathrm{d}t\\
       \geq & \int_0^{\infty} \int_{\{\varphi>\psi+t\}} \theta_{\varphi}^n\,\mathrm{d}t \\
       =& \int_X (\varphi-\psi)\,\theta_{\varphi}^n.
       \end{aligned}
       \]
       The other inequality is similar.
\end{proof}


\begin{lemma}\label{lma:Econcave1}
    The function $E^{\phi}_{\theta}\colon \mathcal{E}^{\infty}(X,\theta;\phi)\rightarrow \mathbb{R}$ is concave.
\end{lemma}
\begin{proof}
    Let $\varphi_0,\varphi_1\in \mathcal{E}^{\infty}(X,\theta;\phi)$. We write $\varphi_t=t\varphi_1+(1-t)\varphi_0$ for each $t\in [0,1]$.

    We claim that for each $t\in [0,1]$, we have
    \begin{equation}\label{eq:ddtEphi_temp1}
        \ddt E^{\phi}_{\theta}(\varphi_t) =\int_X \left(\varphi_1-\varphi_0\right)\,\theta_{\varphi_t}^n.
    \end{equation}
    At the boundary $t=0,1$, the derivative is understood as the one-sided derivative.

    Since $\varphi_0,\varphi_1\in \mathcal{E}^{\infty}(X,\theta;\phi)$, the difference $\varphi_1-\varphi_0$ is quasi-everywhere bounded. Hence all mixed integrals below are finite. We compute the right derivative first. Fix $t\in [0,1)$ and take $s>0$ small enough so that $t+s\leq 1$. By the cocycle formula \eqref{eq:int_by_part_E_temp1},
    \[
        \begin{aligned}
            \frac{E^{\phi}_{\theta}(\varphi_{t+s})-E^{\phi}_{\theta}(\varphi_t)}{s}
            =&\frac{1}{(n+1)s}\sum_{j=0}^n
            \int_X(\varphi_{t+s}-\varphi_t)\,
            \theta_{\varphi_{t+s}}^j\wedge\theta_{\varphi_t}^{n-j}\\
            =&\frac{1}{n+1}\sum_{j=0}^n
            \int_X(\varphi_1-\varphi_0)\,
            \theta_{\varphi_{t+s}}^j\wedge\theta_{\varphi_t}^{n-j}.
        \end{aligned}
    \]
    For each fixed $j$, the last integral is a polynomial in $s$ by the multi-linearity of the non-pluripolar product \itemref{itm:prop:npp1:6}. Letting $s\to0+$, it converges to
    \[
        \int_X(\varphi_1-\varphi_0)\,\theta_{\varphi_t}^n.
    \]
    Thus the right derivative is given by the right-hand side of \eqref{eq:ddtEphi_temp1}. The proof for the left derivative is identical, using $s<0$ with $t+s\in[0,1]$. Hence the left and right derivatives agree for $t\in(0,1)$, and the endpoint derivatives are the corresponding one-sided derivatives.

    Now it suffices to show that the derivative \eqref{eq:ddtEphi_temp1} is decreasing in $t$.
    Take $0\leq s<t \leq 1$. We wish to show that
    \[
        \int_X (\varphi_1-\varphi_0)\,\theta_{\varphi_t}^n\leq \int_X (\varphi_1-\varphi_0)\,\theta_{\varphi_s}^n.
    \]
    This can be equivalently written as
    \[
    \int_X (\varphi_t-\varphi_s)\,\theta_{\varphi_t}^n\leq \int_X (\varphi_t-\varphi_s)\,\theta_{\varphi_s}^n.
    \]
    In particular, replacing $\varphi_0$ by $\varphi_s$ and $\varphi_1$ by $\varphi_t$, we are reduced to proving the following
    \[
        \int_X (\varphi_1-\varphi_0)\,\theta_{\varphi_1}^n\leq \int_X (\varphi_1-\varphi_0)\,\theta_{\varphi_0}^n.
    \]
    But this is just a consequence of \cref{lma:comp_int_ene}.
\end{proof}


\begin{proposition}\label{prop:E_cont_decnet}
    Let $(\varphi_i)_{i\in I}$ be a decreasing net in $\PSH(X,\theta;\phi)$ with
    \[
        \varphi=\inf_{i\in I}\varphi_i\not\equiv -\infty.
    \]
    Then
    \begin{equation}\label{eq:E_cont_decnet}
        \lim_{i\in I} E_{\theta}^{\phi}(\varphi_i)=E_{\theta}^{\phi}(\varphi).
    \end{equation}
\end{proposition}
\begin{proof}
    We first handle the case where $I=\mathbb{Z}_{>0}$. Namely, $(\varphi_i)_i$ is a decreasing sequence.
    Thanks to \cref{lma:EmonoEinfty}, we know that $E_{\theta}^{\phi}(\varphi_i)$ is decreasing in $i$. So the limit in \eqref{eq:E_cont_decnet} exists and the $\geq$ inequality holds. Conversely, let $\psi\in \mathcal{E}^{\infty}(X,\theta;\phi)$ and $\varphi\leq \psi$. We need to show that
    \[
    E_{\theta}^{\phi}(\psi)\geq \lim_{i\to\infty} E_{\theta}^{\phi}\left(\varphi_i\lor \psi\right).
    \]
    Replacing the sequence $(\varphi_i)_i$ by $(\varphi_i\lor\psi)_i$, we are reduced to proving the $\leq$ direction of \eqref{eq:E_cont_decnet} when $\varphi\in \mathcal{E}^{\infty}(X,\theta;\phi)$.

    In this case, thanks to \cref{lma:comp_int_ene}, we have
    \[
    E_{\theta}^{\phi}(\varphi_i)-E_{\theta}^{\phi}(\varphi)\leq \int_X (\varphi_i-\varphi)\,\theta_{\varphi}^n.
    \]
    Our assertion follows from the monotone convergence theorem.

    We now treat the general net case. By \cref{lma:EmonoEinfty}, the net $E_{\theta}^{\phi}(\varphi_i)$ is decreasing and bounded below by $E_{\theta}^{\phi}(\varphi)$. Thus
    \[
        L\coloneqq \lim_{i\in I}E_{\theta}^{\phi}(\varphi_i)
    \]
    exists and satisfies $L\geq E_{\theta}^{\phi}(\varphi)$. It remains to prove the reverse inequality.

    Since $\varphi\leq \varphi_i$ for every $i\in I$, the hypotheses of \cref{cor:decreasingnetL1conv} are satisfied, and hence $\varphi_i\xrightarrow{L^1}\varphi$. For each $k\geq 1$, choose $a_k\in I$ such that
    \[
        \|\varphi_i-\varphi\|_{L^1}<k^{-1},\quad i\geq a_k.
    \]
    We can recursively choose an increasing sequence $(i_k)_{k\geq 1}$ in $I$ such that $i_k\geq a_k$ for all $k$. Then $(\varphi_{i_k})_k$ is a decreasing sequence and $\varphi_{i_k}\to\varphi$ in $L^1$. Its pointwise decreasing limit is therefore equal to $\varphi$ almost everywhere, hence everywhere by \cref{prop:pshfuncdetdense}. Applying the sequence case to $(\varphi_{i_k})_k$, we get
    \[
        \lim_{k\to\infty}E_{\theta}^{\phi}(\varphi_{i_k})=E_{\theta}^{\phi}(\varphi).
    \]
    Since $L\leq E_{\theta}^{\phi}(\varphi_{i_k})$ for every $k$, we conclude $L\leq E_{\theta}^{\phi}(\varphi)$.
\end{proof}

\begin{proposition}\label{prop:E1charac}
    Let $\varphi\in \mathcal{E}(X,\theta;\phi)$. The following are equivalent:
    \begin{enumerate}
        \item\label{itm:prop:E1charac:1} $\varphi\in \mathcal{E}^1(X,\theta;\phi)$;
        \item\label{itm:prop:E1charac:2} $E^{\phi}_{\theta}(\varphi)>-\infty$.
    \end{enumerate}
\end{proposition}
\begin{proof}
    Fix $\varphi\in \mathcal{E}(X,\theta;\phi)$. Without loss of generality, we may assume that $\varphi\leq 0$. In particular, $\varphi\leq P_{\theta}[\varphi]\leq P_{\theta}[\phi]=\phi$.

    \ref{itm:prop:E1charac:2} $\implies$ \ref{itm:prop:E1charac:1}.
    Assume \ref{itm:prop:E1charac:2}. Using the plurifine locality of the non-pluripolar product \itemref{itm:prop:npp1:1}, we have that $\mathds{1}_{\{\varphi>\phi-C\}} \theta_{\varphi\lor (\phi-C)}^n=\mathds{1}_{\{\varphi>\phi-C\}} \theta_{\varphi}^n$ converges \emph{strongly}\footnote{Namely, the convergence holds after integrating against any $L^{\infty}$ function.} to $\theta_{\varphi}^n$ as $C\to \infty$ by the dominated convergence theorem. Therefore, for a fixed $D>0$,
    \[
    \begin{aligned}
    \int_X \Bigl( \varphi\lor (\phi-D)-\phi \Bigr)\,\theta_{\varphi}^n=&\lim_{C\to\infty} \int_{\{\varphi>\phi-C\} } \Bigl( \varphi\lor (\phi-D)-\phi \Bigr)\,\theta_{\varphi\lor (\phi-C)}^n\\
    \geq & \varlimsup_{C\to\infty} \int_{\{\varphi>\phi-C\} } \Bigl( \varphi\lor (\phi-C)-\phi \Bigr)\,\theta_{\varphi\lor (\phi-C)}^n\\
    \geq & \varlimsup_{C\to\infty} \int_X \Bigl( \varphi\lor (\phi-C)-\phi \Bigr)\,\theta_{\varphi\lor (\phi-C)}^n\\
    \geq & \varlimsup_{C\to\infty} (n+1)E_{\theta}^{\phi}\Bigl(\varphi \lor (\phi-C) \Bigr)\\
    \geq & (n+1)E_{\theta}^{\phi}(\varphi).
    \end{aligned}
    \]
    Here the last but one inequality follows from \eqref{eq:Edefbdd} and the fact that $\varphi\lor (\phi-C)\leq \phi$.
    Letting $D\to\infty$, we conclude \ref{itm:prop:E1charac:1}.

    \ref{itm:prop:E1charac:1} $\implies$ \ref{itm:prop:E1charac:2}. Assume \ref{itm:prop:E1charac:1}. Thanks to \cref{lma:comp_int_ene}, we know that for each $C>0$,
    \[
    \begin{aligned}
    E_{\theta}^{\phi}\Bigl(\varphi\lor (\phi-C)\Bigr)\geq & \int_X \Bigl(\varphi\lor (\phi-C)-\phi \Bigr)\,\theta_{\varphi\lor (\phi-C)}^n\\
    =& \int_{\left\{ \varphi>\phi-C \right\}} \Bigl(\varphi-\phi \Bigr)\,\theta_{\varphi \lor (\phi-C)}^n+ \int_{\left\{ \varphi\leq \phi-C \right\}} (-C)\,\theta_{\varphi\lor (\phi-C)}^n \\
    =& \int_{\{\varphi+C>\phi\}} \left(\varphi+C-\phi \right)\,\theta_{\varphi}^n-C\int_X \theta_{\phi}^n\\
    \geq & \int_X \left(\varphi+C-\phi \right)\,\theta_{\varphi}^n-C\int_X \theta_{\phi}^n\\
    =& \int_X (\varphi-\phi)\,\theta_{\varphi}^n.
    \end{aligned}
    \]
    Due to \cref{prop:E_cont_decnet}, we have
    \[
    \lim_{C\to\infty}E_{\theta}^{\phi}\Bigl(\varphi\lor (\phi-C)\Bigr)=E_{\theta}^{\phi}(\varphi),
    \]
    so \ref{itm:prop:E1charac:2} holds.
\end{proof}

\begin{corollary}\label{cor:E_concave2}
    The function $E^{\phi}_{\theta}\colon \PSH(X,\theta;\phi)\rightarrow [-\infty,\infty)$ is concave.
\end{corollary}
\begin{proof}
    Let $\varphi,\psi\in \PSH(X,\theta;\phi)$ and $t\in (0,1)$. We need to show that
    \begin{equation}\label{eq:Econcave_temp1}
    E^{\phi}_{\theta}\Bigl( t\varphi+(1-t)\psi \Bigr)\geq tE^{\phi}_{\theta}(\varphi)+(1-t)E^{\phi}_{\theta}(\psi).
    \end{equation}

    For this purpose, observe that when $C\to \infty$, we have
    \[
        t\Bigl(\varphi\lor (\phi-C)\Bigr)+(1-t)\Bigl(\psi\lor (\phi-C)\Bigr)\searrow t\varphi+(1-t)\psi.
    \]
    Observe that
    \[
        \varphi\lor (\phi-C),\psi\lor (\phi-C)\in \mathcal{E}^{\infty}(X,\theta;\phi).
    \]
    By \cref{lma:Econcave1}, for each $C>0$, we have
    \[
        E_{\theta}^{\phi}\Biggl( t\Bigl(\varphi\lor (\phi-C)\Bigr)+(1-t)\Bigl(\psi\lor (\phi-C)\Bigr) \Biggr)\geq tE_{\theta}^{\phi}\Bigl(\varphi\lor (\phi-C)\Bigr)+(1-t)E_{\theta}^{\phi}\Bigl(\psi\lor (\phi-C)\Bigr).
    \]
    Letting $C\to \infty$, we conclude \eqref{eq:Econcave_temp1} using \cref{prop:E_cont_decnet}.
\end{proof}

\begin{corollary}\label{cor:E_classes_convex}
    The sets $\PSH(X,\theta;\phi)$, $\mathcal{E}(X,\theta;\phi)$, $\mathcal{E}^1(X,\theta;\phi)$, $\mathcal{E}^{\infty}(X,\theta;\phi)$ are all convex.
\end{corollary}
\begin{proof}
    The sets $\PSH(X,\theta;\phi)$ and $\mathcal{E}^{\infty}(X,\theta;\phi)$ are trivially convex.

    We next handle $\mathcal{E}(X,\theta;\phi)$. In this case, suppose that $\varphi,\psi \in \mathcal{E}(X,\theta;\phi)$ and $t\in (0,1)$. Then we compute
    \[
    \begin{split}
        \int_X \theta_{t\varphi+(1-t)\psi}^n=\sum_{j=0}^n \binom{n}{j}t^j (1-t)^{n-j} \int_X \theta_{\varphi}^j\wedge \theta_{\psi}^{n-j}\\
        =\sum_{j=0}^n \binom{n}{j}t^j (1-t)^{n-j} \int_X \theta_{\phi}^n=\int_X \theta_{\phi}^n,
    \end{split}
    \]
    where we applied \cref{prop:Ppresmass} in the second equality. Therefore, $t\varphi+(1-t)\psi \in \mathcal{E}(X,\theta;\phi)$.

    Finally, let us consider $\mathcal{E}^1(X,\theta;\phi)$. In this case, suppose that $\varphi,\psi \in \mathcal{E}^1(X,\theta;\phi)$ and $t\in (0,1)$. We have already shown that $t\varphi+(1-t)\psi \in \mathcal{E}(X,\theta;\phi)$. Furthermore, the energy of this potential is finite due to \cref{cor:E_concave2}. Therefore, $t\varphi+(1-t)\psi \in \mathcal{E}^1(X,\theta;\phi)$ in view of \cref{prop:E1charac}.
\end{proof}


\begin{lemma}\label{lma:intdom1}
	Let $\varphi,\psi,\gamma\in \mathcal{E}(X,\theta;\phi)$. Assume that $\gamma\geq \varphi\lor\psi$.
	Then
	\[
	\int_X (\gamma-\varphi)\,\theta_{\psi}^n\leq 2\int_X (\gamma-\varphi)\,\theta_{\varphi}^n+2\int_X (\gamma-\psi)\,\theta_{\psi}^n.
	\]
\end{lemma}
\begin{proof}
	Observe that
	\begin{equation}\label{eq:gvp2}
		\{\gamma>\varphi+2t\}\subseteq \{\gamma>\psi+t \}\cup \{\psi>\varphi+t\}\,.
	\end{equation}
	So
	\[
	\begin{aligned}
		\int_X (\gamma-\varphi)\,\theta_{\psi}^n
		=& 2\int_0^{\infty} \int_{\{\gamma>\varphi+2t\}}\,\theta_{\psi}^n\,\mathrm{d}t \\
		\leq& 2\int_0^{\infty} \int_{\{\gamma>\psi+t\}}\,\theta_{\psi}^n\,\mathrm{d}t+
		2\int_0^{\infty} \int_{\{\psi>\varphi+t\}}\,\theta_{\psi}^n\,\mathrm{d}t\\
		\leq & 2\int_X (\gamma-\psi) \,\theta_{\psi}^n+2\int_0^{\infty} \int_{\{\psi>\varphi+t\}}\,\theta_{\varphi}^n\,\mathrm{d}t\\
		\leq & 2\int_X (\gamma-\psi) \,\theta_{\psi}^n+2\int_0^{\infty} \int_{\{\gamma>\varphi+t\}}\,\theta_{\varphi}^n\,\mathrm{d}t\\
		=&2\int_X (\gamma-\psi) \,\theta_{\psi}^n+2\int_X (\gamma-\varphi) \,\theta_{\varphi}^n,
	\end{aligned}
	\]
    where the second line follows from \eqref{eq:gvp2}, while the third line follows from \cref{prop:compa_princ}.
\end{proof}

\begin{corollary}\label{cor:mix_MA_energy_finite}
    Let $\varphi,\psi,\gamma_1,\ldots,\gamma_n\in \mathcal{E}^1(X,\theta;\phi)$. Then
    \begin{equation}\label{eq:mix_MA_energy}
        \int_X (\psi-\varphi)\,\theta_{\gamma_1}\wedge \cdots \wedge \theta_{\gamma_n}
    \end{equation}
    is well-defined and is in $\mathbb{R}$.
\end{corollary}
\begin{proof}
    \textbf{Step~1}. We first consider the case $\psi=\phi$.

    We can find $C>0$ so that $\varphi\leq \phi+C$. Therefore, $\phi-\varphi \geq -C$ almost everywhere with respect to $\theta_{\gamma_1}\wedge \cdots \wedge \theta_{\gamma_n}$. So the integral \eqref{eq:mix_MA_energy} is well-defined and takes value in $\mathbb{R}\cup \{\infty\}$. We show that the integral is finite. For this purpose, we may replace $\varphi$ by $\varphi-C$ and assume that $\varphi\leq \phi$. Put
    \[
        \eta=n^{-1}(\gamma_1+\cdots+\gamma_n).
    \]
    Then by the multi-linearity of the non-pluripolar product \itemref{itm:prop:npp1:6} and positivity, we have
    \[
        \int_X (\phi-\varphi)\,\theta_{\gamma_1}\wedge \cdots \wedge \theta_{\gamma_n}\leq n^n\int_X (\phi-\varphi)\,\theta_{\eta}^n.
    \]
    Thanks to \cref{cor:E_classes_convex}, $\eta\in \mathcal{E}^1(X,\theta;\phi)$. Choose $A>0$ so that $\eta\leq \phi+A$. Applying \cref{lma:intdom1} to $\varphi,\eta$ and $\phi+A$, we get
    \[
        \int_X(\phi+A-\varphi)\,\theta_{\eta}^n
        \leq 2\int_X(\phi+A-\varphi)\,\theta_{\varphi}^n
        +2\int_X(\phi+A-\eta)\,\theta_{\eta}^n<\infty.
    \]
    Hence $\int_X(\phi-\varphi)\,\theta_{\eta}^n<\infty$, and Step~1 follows.

    \textbf{Step~2}. We handle the general case. Take $C>0$ so that
    \[
        \psi\leq \phi+C,\quad \varphi\leq \phi+C.
    \]
    Then
    \[
        (\psi-C)-\phi \leq \psi-\varphi\leq \phi-(\varphi-C).
    \]
    Step~1 applied to $\varphi-C$ and to $\psi-C$ shows that $\phi-(\varphi-C)$ and $\phi-(\psi-C)$ are integrable with respect to $\theta_{\gamma_1}\wedge \cdots \wedge \theta_{\gamma_n}$. Hence the positive and negative parts of $\psi-\varphi$ are both integrable, so \eqref{eq:mix_MA_energy} is well-defined and finite.
\end{proof}

\begin{proposition}\label{prop:cocycE1}
    Given $\varphi,\psi\in \mathcal{E}^1(X,\theta;\phi)$, we have the following cocycle equality
    \begin{equation}\label{eq:Ecocyc}
        E^{\phi}_{\theta}(\psi)-E^{\phi}_{\theta}(\varphi)=\frac{1}{n+1}\sum_{j=0}^n \int_X (\psi-\varphi)\, \theta_{\psi}^j\wedge \theta_{\varphi}^{n-j}.
    \end{equation}
\end{proposition}
As a consequence of \eqref{eq:Ecocyc}, for $\varphi\in \mathcal{E}^1(X,\theta;\phi)$, \eqref{eq:Edefbdd} continues to hold. Note that each term on the right-hand side of \eqref{eq:Ecocyc} is finite due to \cref{cor:mix_MA_energy_finite}.
\begin{proof}
    Since the conclusion is known when $\varphi,\psi\in \mathcal{E}^{\infty}(X,\theta;\phi)$ (see \eqref{eq:int_by_part_E_temp1}), it suffices to prove that for each $k=0,\ldots,n$, we have
    \begin{equation}\label{eq:canonic_app_energy_conv}
    \lim_{C\to\infty} \int_X \Bigl( \varphi\lor (\phi-C)-\psi\lor (\phi-C) \Bigr) \,\theta_{\varphi\lor (\phi-C)}^k\wedge \theta_{\psi\lor (\phi-C)}^{n-k}=\int_X (\varphi-\psi)\,\theta_{\varphi}^k\wedge \theta_{\psi}^{n-k}.
    \end{equation}
    For this purpose, we may assume that $\varphi,\psi\leq \phi$ and it suffices to establish the following:
    \begin{equation}\label{eq:intvarphilorphimCmphi_temp1}
        \lim_{C\to\infty} \int_X \Bigl( \varphi\lor (\phi-C)-\phi \Bigr) \,\theta_{\varphi\lor (\phi-C)}^k\wedge \theta_{\psi\lor (\phi-C)}^{n-k}=\int_X (\varphi-\phi)\,\theta_{\varphi}^k\wedge \theta_{\psi}^{n-k}.
    \end{equation}
    Fix $C>0$. We compute the difference as follows:
    \[
    \begin{aligned}
    &\int_X \Bigl( \varphi\lor (\phi-C)-\phi \Bigr) \,\theta_{\varphi\lor (\phi-C)}^k\wedge \theta_{\psi\lor (\phi-C)}^{n-k}-\int_X (\varphi-\phi)\,\theta_{\varphi}^k\wedge \theta_{\psi}^{n-k}\\
    =& \int_{\{\min\{\varphi,\psi\}>\phi-C\}}\left( \varphi-\phi \right) \,\theta_{\varphi}^k\wedge \theta_{\psi}^{n-k}-\int_X (\varphi-\phi)\,\theta_{\varphi}^k\wedge \theta_{\psi}^{n-k}\\
    &+\int_{\{\min\{\varphi,\psi\}\leq \phi-C\}}\Bigl( \varphi\lor (\phi-C)-\phi \Bigr) \,\theta_{\varphi\lor (\phi-C)}^k\wedge \theta_{\psi\lor (\phi-C)}^{n-k}\\
    =& -\int_{\{\min\{\varphi,\psi\}\leq \phi-C\}}\left( \varphi-\phi \right) \,\theta_{\varphi}^k\wedge \theta_{\psi}^{n-k}+\\
     &+\int_{\{\min\{\varphi,\psi\}\leq \phi-C\}}\Bigl( \varphi\lor (\phi-C)-\phi \Bigr) \,\theta_{\varphi\lor (\phi-C)}^k\wedge \theta_{\psi\lor (\phi-C)}^{n-k}
    \end{aligned}
    \]
    We first handle the first term. In order to show that it converges to $0$ when $C\to \infty$, it suffices to show that
    \[
    \int_X\left( \phi-\varphi \right) \,\theta_{\varphi}^k\wedge \theta_{\psi}^{n-k}<\infty.
    \]
    For this purpose, it suffices to prove a stronger inequality:
    \begin{equation}\label{eq:intXphimvarphifinite_temp1}
    \int_X\left( \phi-\varphi \right) \,\theta_{(\varphi+\psi)/2}^n<\infty.
    \end{equation}
    From \cref{cor:E_classes_convex}, we know that $(\varphi+\psi)/2\in \mathcal{E}^1(X,\theta;\phi)$. Therefore, \eqref{eq:intXphimvarphifinite_temp1} follows from \cref{lma:intdom1}.


    It remains to handle the second term. It then suffices to prove the following:
    \begin{equation}\label{eq:limCCintvarphiphiC_temp1}
    \begin{split}
    \lim_{C\to\infty} C\int_{\{\varphi\leq \phi-C\}} \theta_{\varphi\lor (\phi-C)}^k\wedge \theta_{\psi\lor (\phi-C)}^{n-k}=&0,\\
    \lim_{C\to\infty} C\int_{\{\psi\leq \phi-C\}} \theta_{\varphi\lor (\phi-C)}^k\wedge \theta_{\psi\lor (\phi-C)}^{n-k}=&0.
    \end{split}
    \end{equation}
    By symmetry, it suffices to handle the first. Observe the following inclusions:
    \[
    \{\varphi\leq \phi-C\}\subseteq \left\{ \varphi\lor (\phi-C)\leq \frac{1}{2}\Bigl( \psi\lor (\phi-C)+\phi-C \Bigr) \right\}\subseteq \{\varphi\leq \phi-C/2\}.
    \]
    By \cref{prop:compa_princ}, we find
    \[
    \begin{aligned}
        &C\int_{\{\varphi\leq \phi-C\}} \theta_{\varphi\lor (\phi-C)}^k\wedge \theta_{\psi\lor (\phi-C)}^{n-k}\\
        \leq & C\int_{\left\{\varphi\lor (\phi-C)\leq \frac{1}{2}\bigl(\psi\lor (\phi-C)+\phi-C\bigr)\right\}} \theta_{\varphi\lor (\phi-C)}^k\wedge \theta_{\psi\lor (\phi-C)}^{n-k}\\
        \leq & 2^{n-k}C\int_{\left\{\varphi\lor (\phi-C)\leq \frac{1}{2}\bigl(\psi\lor (\phi-C)+\phi-C\bigr)\right\}} \theta_{\varphi\lor (\phi-C)}^k\wedge \theta_{\frac{1}{2}\bigl(\psi\lor (\phi-C)+\phi-C\bigr)}^{n-k}\\
        \leq & 2^{n-k}C\int_{\{\varphi\lor (\phi-C)\leq \frac{1}{2}\bigl(\psi\lor (\phi-C)+\phi-C\bigr)\}} \theta_{\varphi\lor (\phi-C)}^n\\
        \leq & 2^{n-k}C\int_{\{\varphi\leq \phi-C/2\}} \theta_{\varphi\lor (\phi-C)}^n\\
        = & 2^{n-k}C\int_{\{\varphi\leq \phi-C/2\}} \theta_{\varphi}^n.
    \end{aligned}
    \]
    Here on the final step, we have applied the following argument:
    \[
    \begin{aligned}
    &\int_{\left\{\varphi\leq \phi-C/2 \right\}} \theta_{\varphi\lor (\phi-C)}^n &\\
    = &\int_X \theta_{\varphi\lor (\phi-C)}^n-\int_{\left\{\varphi> \phi-C/2 \right\}} \theta_{\varphi\lor (\phi-C)}^n &\\
    = & \int_X \theta_{\varphi}^n-\int_{\left\{\varphi> \phi-C/2 \right\}} \theta_{\varphi}^n & \textup{by \itemref{itm:prop:npp1:1}}\\
    =& \int_{\left\{\varphi\leq \phi-C/2 \right\}} \theta_{\varphi}^n. &
    \end{aligned}
    \]

    On the other hand,
    \[
    \varlimsup_{C\to\infty} C\int_{\{\varphi\leq \phi-C/2 \}} \theta_{\varphi}^n\leq 2 \varlimsup_{C\to\infty} \int_{\{\varphi\leq \phi-C \}}(\phi-\varphi)\, \theta_{\varphi}^n=0,
    \]
    since $\varphi\in \mathcal{E}^1(X,\theta;\phi)$. We conclude \eqref{eq:limCCintvarphiphiC_temp1}.
\end{proof}
\begin{proposition}\label{prop:Ediffbdd}
Let $\varphi,\psi\in \mathcal{E}^1(X,\theta;\phi)$. Then
    \begin{equation}
        \int_X \left(\psi-\varphi\right)\theta^n_{\psi}\leq E_{\theta}^{\phi}(\psi)-E_{\theta}^{\phi}(\varphi)\leq \int_X \left(\psi-\varphi\right)\theta^n_{\varphi}.
    \end{equation}
\end{proposition}
\begin{proof}
    Thanks to \eqref{eq:canonic_app_energy_conv}, this can be reduced to the case where $\varphi,\psi\in \mathcal{E}^{\infty}(X,\theta;\phi)$, which is established in \cref{lma:comp_int_ene}.
\end{proof}


\begin{proposition}\label{prop:Econt_mono}
    Let $(\varphi_i)_{i\in I}$ be a net in $\mathcal{E}^1(X,\theta;\phi)$. Assume that $\varphi\in \mathcal{E}^1(X,\theta;\phi)$ and either one of the following conditions holds:
    \begin{enumerate}
        \item\label{itm:prop:Econt_mono:1} $(\varphi_i)_i$ is a decreasing net with $\varphi=\inf_{i\in I}\varphi_i$;
        \item\label{itm:prop:Econt_mono:2} $(\varphi_i)_i$ is an increasing net with almost everywhere limit $\varphi$.
    \end{enumerate}
    Then
    \begin{equation}\label{eq:Iconv_mono}
    \lim_{i\in I} E_{\theta}^{\phi}(\varphi_i)=E_{\theta}^{\phi}(\varphi).
    \end{equation}
\end{proposition}
\begin{proof}
    The decreasing case is already proved in \cref{prop:E_cont_decnet}, so let us focus on the case where $(\varphi_i)_i$ is increasing.

    We first handle the case where $I=\mathbb{Z}_{>0}$. Namely, $(\varphi_i)_i$ is an increasing sequence.

    \textbf{Step~1}. We first prove the result when $\varphi_j,\varphi\in \mathcal{E}^{\infty}(X,\theta;\phi)$ for each $j\geq 1$.

    Due to \cref{cor:incseqnppcont}, we have
    \[
        \lim_{j\to\infty} \int_X \theta_{\varphi}^k\wedge \theta_{\varphi_j}^{n-k}=\int_X \theta_{\varphi}^n
    \]
    for each $k=0,\ldots,n$.
    By \eqref{eq:int_by_part_E_temp1}, we have
    \[
    \begin{aligned}
    0\leq & E_{\theta}^{\phi}(\varphi)-E_{\theta}^{\phi}(\varphi_j)\\
    =& \frac{1}{n+1}\sum_{k=0}^n \int_X(\varphi-\varphi_j)\,\theta_{\varphi}^k\wedge \theta_{\varphi_j}^{n-k}.
    \end{aligned}
    \]
    We then apply \cref{thm:semicon} to conclude that the right-hand side converges to $0$.

    \textbf{Step~2}. We prove the general case. The $\leq$ direction in \eqref{eq:Iconv_mono} follows from the monotonicity of $E_{\theta}^{\phi}$, as proved in \cref{lma:EmonoEinfty}. It remains to prove the $\geq$ direction.

    After subtracting a constant from all $\varphi_j$ and from $\varphi$, which shifts all the energies by the same constant in view of \eqref{eq:Ephiaddcst}, we may assume that $\varphi\leq \phi$. For each $C>0$ and $j\geq 1$, we let
    \[
    \varphi_C\coloneqq \varphi\lor (\phi-C),\quad \varphi_{j,C}\coloneqq \varphi_j\lor (\phi-C).
    \]
    By Step~1, for each $C>0$, we have
    \[
    \lim_{j\to\infty} E_{\theta}^{\phi}(\varphi_{j,C})= E_{\theta}^{\phi}(\varphi_{C})\geq E_{\theta}^{\phi}(\varphi).
    \]
    It suffices therefore to show that
    \begin{equation}\label{eq:limCCEphi_temp1}
    \lim_{C\to \infty} E_{\theta}^{\phi}(\varphi_{j,C})-E_{\theta}^{\phi}(\varphi_{j})=0\quad \textup{uniformly for $j$ in a tail}.
    \end{equation}
    Fix $i\geq 1$.
    For $j\geq i$, we compute
    \[
    \begin{aligned}
        & E_{\theta}^{\phi}(\varphi_{j,C})-E_{\theta}^{\phi}(\varphi_{j})\\
        \leq & \int_X \left( \varphi_{j,C}-\varphi_{j}\right)\,\theta_{\varphi_{j}}^n &\textup{by \cref{prop:Ediffbdd}}\\
        \leq & \int_{\{\varphi_j\leq \phi-C\}} (\phi-C-\varphi_j)\,\theta_{\varphi_j}^n\\
        =& \int_C^{\infty}\int_{\{\varphi_j\leq \phi-t\}} \theta_{\varphi_j}^n\,\mathrm{d}t\\
        \leq & \int_C^{\infty}\int_{\{\varphi_i\leq (\varphi_j+\phi-t)/2\}} \theta_{\varphi_j}^n\,\mathrm{d}t\\
        \leq & 2^n\int_C^{\infty}\int_{\{\varphi_i\leq (\varphi_j+\phi-t)/2\}} \theta_{(\varphi_j+\phi-t)/2}^n\,\mathrm{d}t\\
        \leq & 2^n\int_C^{\infty}\int_{\{\varphi_i\leq (\varphi_j+\phi-t)/2\}} \theta_{\varphi_i}^n\,\mathrm{d}t & \textup{by \cref{prop:compa_princ}}\\
        \leq &2^n\int_C^{\infty}\int_{\{\varphi_i\leq \phi-t/2\}} \theta_{\varphi_i}^n\,\mathrm{d}t\\
        = & 2^{n+1}\int_{\{\varphi_i< \phi-C/2\}} (\phi-\varphi_i-C/2)\,\theta_{\varphi_i}^n\\
        = & 2^{n+1}\int_{\{\varphi_i< \phi-C/2\}} (\phi-\varphi_i)\,\theta_{\varphi_i}^n-C 2^n\int_{\{\varphi_i< \phi-C/2\}}\theta_{\varphi_i}^n.
    \end{aligned}
    \]
    As $C\to \infty$, the first term converges to $0$ since $\varphi_i\in \mathcal{E}^1(X,\theta;\phi)$. As for the second term,
    \[
    \lim_{C\to\infty}C\int_{\{\varphi_i< \phi-C/2 \}} \theta_{\varphi_i}^n\leq 2 \lim_{C\to\infty}\int_{\{\varphi_i< \phi-C/2 \}}(\phi-\varphi_i)\, \theta_{\varphi_i}^n = 0.
    \]
    We conclude \eqref{eq:limCCEphi_temp1}.

    We now pass from sequences to nets. By \cref{lma:EmonoEinfty}, the net $(E_{\theta}^{\phi}(\varphi_i))_{i\in I}$ is increasing and bounded above by $E_{\theta}^{\phi}(\varphi)$. Let
    \[
        L\coloneqq \lim_{i\in I}E_{\theta}^{\phi}(\varphi_i).
    \]
    Then $L\leq E_{\theta}^{\phi}(\varphi)$.

    By \cref{prop:pshfuncdetdense}, we have
    \[
        \varphi=\uscsup[i\in I]\varphi_i.
    \]
    Moreover, $\varphi_i\leq \varphi$ for all $i\in I$, so the family is locally uniformly bounded above.
    Applying Choquet's lemma \cref{prop:Choquet}, we can find a countable subset $J\subseteq I$ such that
    \[
        \varphi=\uscsup[j\in J]\varphi_j.
    \]
    Enumerate $J=\{j_1,j_2,\ldots\}$, repeating elements if necessary. Since $I$ is directed, we can choose an increasing sequence $(i_k)_{k\geq 1}$ in $I$ such that $i_k$ dominates $j_1,\ldots,j_k$. It follows that $(\varphi_{i_k})_k$ increases almost everywhere to $\varphi$. Applying the sequence case to this sequence gives
    \[
        \lim_{k\to\infty}E_{\theta}^{\phi}(\varphi_{i_k})=E_{\theta}^{\phi}(\varphi).
    \]
    Since $E_{\theta}^{\phi}(\varphi_{i_k})\leq L$ for every $k$, we conclude $E_{\theta}^{\phi}(\varphi)\leq L$.
\end{proof}

Next we want to prove that $\mathcal{E}^1(X,\theta;\phi)$ is closed under the rooftop operator. For this purpose, we shall need a few preliminary results.

We shall approximate the rooftop operator by solutions of certain Monge--Ampère equations.
\begin{lemma}\label{lma:rooftop_MAequation}
	Let $\varphi,\psi\in \mathcal{E}^{\infty}(X,\theta;\phi)$. Then there is $\gamma\in \mathcal{E}^{\infty}(X,\theta;\phi)$ such that
	\begin{equation}\label{eq:solution_thetavarphiephiu}
		\theta_{\gamma}^n=\mathrm{e}^{\gamma-\varphi}\,\theta_{\varphi}^n+\mathrm{e}^{\gamma-\psi}\,\theta_{\psi}^n.
	\end{equation}
\end{lemma}
It is not clear if $\int_X\mathrm{e}^{-\varphi}\theta_{\varphi}^n<\infty$, hence we cannot say that $\mathrm{e}^{-\varphi}\theta_{\varphi}^n+\mathrm{e}^{-\psi}\theta_{\psi}^n$ is a Radon measure, so \cref{thm:AY_rel} is not directly applicable.
\begin{proof}
	For each $j\geq 1$, let $\varphi_j\coloneqq \varphi\lor(-j)$, $\psi_j\coloneqq \psi\lor(-j)$. These functions are not necessarily $\theta$-psh.
	Let
	\[
	\mu_j=\mathrm{e}^{-\varphi_j}\theta_{\varphi}^n+\mathrm{e}^{-\psi_j}\theta_{\psi}^n.
	\]
	By \cref{thm:AY_rel}, we can find $\gamma_j\in \mathcal{E}(X,\theta;\phi)$ such that
	\begin{equation}\label{eq:thetagammaj_temp1}
		\theta_{\gamma_j}^n=\mathrm{e}^{\gamma_j}\mu_j=\mathrm{e}^{\gamma_j-\varphi_j}\theta_{\varphi}^n+\mathrm{e}^{\gamma_j-\psi_j}\theta_{\psi}^n.
	\end{equation}
	Take a constant $C>0$ so that $\psi-2C\leq \varphi\leq \psi+2C$.
	Let
	\[
	\eta\coloneqq \frac{\varphi+\psi}{2}-C-n\log 2.
	\]
	Then $\eta\in \mathcal{E}^{\infty}(X,\theta;\phi)$ and we have the following inequalities:
    \[
    \begin{aligned}
    \theta_{\eta}^n =& \theta_{(\varphi+\psi)/2}^n\\
    \geq & 2^{-n}\theta_{\varphi}^n+2^{-n}\theta_{\psi}^n\\
    \geq & 2^{-n}\mathrm{e}^{(\varphi+\psi)/2}\mathrm{e}^{-C}\left( \mathrm{e}^{-\varphi}\theta_{\varphi}^n+\mathrm{e}^{-\psi}\theta_{\psi}^n \right)\\
    \geq & 2^{-n}\mathrm{e}^{(\varphi+\psi)/2}\mathrm{e}^{-C}\left( \mathrm{e}^{-\varphi_j}\theta_{\varphi}^n+\mathrm{e}^{-\psi_j}\theta_{\psi}^n \right)\\
    = &
    \mathrm{e}^{\eta}\mu_j.
    \end{aligned}
    \]
	Hence, $\gamma_j\geq \eta$ by \cref{thm:AY_rel}. By \cref{thm:AY_rel} again, $\gamma_j$ is decreasing in $j$; let $\gamma=\inf_{j>0}\gamma_j$.
	Then $\gamma\geq \eta$, hence $\gamma\in \mathcal{E}^{\infty}(X,\theta;\phi)$.

    Now observe that
    \[
    \theta_{\gamma_j}^n\rightharpoonup \theta_{\gamma}^n
    \]
    as $j\to \infty$, by \cref{thm:semicon}.
    Finally observe that there is a constant $C'>0$ so that
    \[
    \gamma_1\leq \varphi+C',\quad \gamma_1\leq \psi+C',
    \]
    since $\varphi,\psi\in \mathcal{E}^{\infty}(X,\theta;\phi)$.
    Therefore, for each $j>0$,
    \[
    \gamma_j\leq \varphi_j+C',\quad \gamma_j\leq \psi_j+C'.
    \]
    Hence, \eqref{eq:solution_thetavarphiephiu} follows from \eqref{eq:thetagammaj_temp1} and the dominated convergence theorem by letting $j\to\infty$.
\end{proof}


We also need a few integral estimates.


\begin{lemma}\label{lma:fundam_estimate}
	Let $\varphi,\psi,\gamma\in \mathcal{E}(X,\theta;\phi)$. Assume that $\varphi\leq \psi\leq \gamma$. Then
	\[
	\int_X (\gamma-\psi)\,\theta_{\psi}^n\leq 2^{n+1}\int_X (\gamma-\varphi)\,\theta_{\varphi}^n.
	\]
\end{lemma}
\begin{proof}
	Observe that for any $t\geq 0$,
	\begin{equation}\label{eq:ppgcomp}
        \begin{split}
		\{\gamma>\psi+2t\}\setminus \{\psi=-\infty\}\subseteq \left\{(\gamma+\psi)/2>\psi+t\right\}\\
        \subseteq \left\{(\gamma+\psi)/2>\varphi+t\right\}\subseteq \{\gamma>\varphi+t\}.
        \end{split}
	\end{equation}
	So
	\[
	\begin{aligned}
		\int_X (\gamma-\psi)\,\theta_{\psi}^n
		=& 2\int_0^{\infty}  \int_{\{\gamma-\psi>2t\}}\,\theta_{\psi}^n\,\mathrm{d}t &\\
		\leq & 2\int_0^{\infty} \int_{\left\{(\gamma+\psi)/2>\psi+t\right\}}\,\theta_{\psi}^n\,\mathrm{d}t & \textup{by \eqref{eq:ppgcomp}}\\
		\leq  & 2^{n+1}\int_0^{\infty} \int_{\left\{(\gamma+\psi)/2>\varphi+t\right\}}\,\theta_{(\gamma+\psi)/2}^n\,\mathrm{d}t &\\
		\leq & 2^{n+1}\int_0^{\infty} \int_{\left\{(\gamma+\psi)/2>\varphi+t\right\}}\,\theta_{\varphi}^n\,\mathrm{d}t & \textup{by \cref{prop:compa_princ}}\\
		\leq &2^{n+1}\int_0^{\infty} \int_{\left\{\gamma>\varphi+t\right\}}\,\theta_{\varphi}^n\,\mathrm{d}t &\textup{by \eqref{eq:ppgcomp}}\\
		=& 2^{n+1} \int_X (\gamma-\varphi)\,\theta_{\varphi}^n &.
	\end{aligned}
	\]
	This completes the proof.
\end{proof}

\begin{lemma}\label{lma:stab}
	Let $\varphi_j\in \mathcal{E}(X,\theta;\phi)$ ($j\in \mathbb{Z}_{>0}$) and $\gamma\in \mathcal{E}(X,\theta;\phi)$. Assume that $\varphi_j\leq \gamma$ for each $j$ and that $\varphi_j$ converges to $\varphi\in \PSH(X,\theta)$ with respect to the $L^1$-topology.
	Assume that there is a constant $A>0$ such that for any $j>0$,
	\[
	\int_X (\varphi_j-\gamma)\,\theta_{\varphi_j}^n\geq -A.
	\]
	Then $\varphi\in \mathcal{E}(X,\theta;\phi)$ and
	\begin{equation}\label{eq:intvarphigammageqm2np3A}
	    \int_X (\varphi-\gamma)\,\theta_{\varphi}^n\geq - 2^{n+3}A.
	\end{equation}
\end{lemma}
\begin{proof}
    Replacing simultaneously $\gamma,\varphi_j,\varphi$ by $\gamma-B,\varphi_j-B,\varphi-B$ for $B\gg 1$, we may assume that $\gamma\leq 0$. The assumptions and the conclusion are invariant under this translation. Since $\gamma\in\mathcal E(X,\theta;\phi)$ and $P_\theta[\phi]=\phi$, we then have $\gamma\leq\phi$.

	\textbf{Step~1}. Assume that $(\varphi_j)_j$ is decreasing. In this case, we prove
	\begin{equation}\label{eq:intvarphigammageqm4A_temp1}
	\int_X (\varphi-\gamma)\,\theta_{\varphi}^n\geq -4A.
    \end{equation}
	By \cref{lma:intdom1}, for any $j, k>0$,
	\[
	\int_X (\varphi_j-\gamma)\,\theta_{\varphi_k}^n\geq -4A.
	\]
	For any $C>0$,
	\[
	\int_X \Bigl(\varphi_j \lor (\gamma-C)-\gamma\Bigr)\,\theta_{\varphi_k}^n\geq \int_X (\varphi_j-\gamma) \,\theta_{\varphi_k}^n\geq -4A.
	\]
	Letting $k\to \infty$, by \cref{thm:semicon}, we find
	\[
	\int_X \Bigl(\varphi_j \lor (\gamma-C)-\gamma\Bigr)\,\theta_{\varphi}^n\geq -4A.
	\]
	Letting $j\to \infty$, by the monotone convergence theorem, we get
	\[
	\int_X \Bigl(\varphi \lor (\gamma-C)-\gamma\Bigr)\,\theta_{\varphi}^n\geq -4A.
	\]
	Then we let $C\to \infty$, again by the monotone convergence theorem, we conclude \eqref{eq:intvarphigammageqm4A_temp1}.

	\textbf{Step~2}. In general, let
    \[
    \psi_j=\uscsup[k\geq j] \varphi_k.
    \]
    Then $\varphi=\inf_{j>0}\psi_j$ due to \cref{cor:L1limipp}.

	For each $C>0$, let
	\[
	\psi_{j,C}=\psi_j \lor (\gamma-C),\quad \varphi_C=\varphi \lor (\gamma-C).
	\]
    Then $\psi_{j,C},\varphi_C\in \mathcal{E}(X,\theta;\phi)$ since $\gamma-C\in \mathcal{E}(X,\theta;\phi)$ thanks to the monotonicity theorem \cref{thm:mono}.
	Observe that $\psi_{j,C}$ decreases to $\varphi_C$ as $j\to \infty$. Moreover,
    \[
    \gamma\geq \psi_{j,C} \geq \psi_j\geq \varphi_j.
    \]
	By \cref{lma:fundam_estimate},
	\[
	\int_X (\psi_{j,C}-\gamma)\,\theta_{\psi_{j,C}}^n\geq -2^{n+1}A.
	\]
	By Step~1,
	\begin{equation}\label{eq:inproof1}
		\int_X(\varphi_C-\gamma)\,\theta_{\varphi_C}^n\geq -2^{n+3}A.
	\end{equation}
	In particular, by the plurifine locality of the non-pluripolar product \itemref{itm:prop:npp1:1},
	\[
	    \int_{\{\varphi>\gamma-C\}}(\varphi-\gamma)\,\theta_{\varphi}^n\geq -2^{n+3}A.
	\]
	Letting $C\to \infty$, we conclude \eqref{eq:intvarphigammageqm2np3A}. As a consequence,
    \[
        \int_X (\gamma-\varphi)\,\theta_{\varphi}^n<\infty.
    \]
	We still have to prove that $\varphi\in \mathcal{E}(X,\theta;\phi)$. In fact, by \eqref{eq:inproof1}, for any $C>0$,
	\[
	\int_{\{\varphi\leq \gamma-C\}}\,\theta_{\varphi_C}^n\leq \frac{1}{C}\int_X (\gamma-\varphi_C)\,\theta_{\varphi_C}^n\leq 2^{n+3}C^{-1}A.
	\]
	Using the monotonicity theorem \cref{thm:mono} and the plurifine locality of the non-pluripolar product \itemref{itm:prop:npp1:1}, we find
    \[
    \int_X \theta_{\phi}^n=\int_X \theta_{\varphi_C}^n=\int_{\{\varphi\leq \gamma-C\}}\,\theta_{\varphi_C}^n+\int_{\{\varphi>\gamma-C\}}\,\theta_{\varphi}^n.
    \]
    Letting $C\to\infty$, we conclude that
    \[
    \int_X \theta_{\varphi}^n=\int_X \theta_{\phi}^n.
    \]
    This completes the proof.
\end{proof}

\begin{proposition}\label{prop:relrooftopclosed}
    Assume that $\varphi,\psi\in \mathcal{E}(X,\theta;\phi)$ (resp. $\mathcal{E}^1(X,\theta;\phi)$, $\mathcal{E}^{\infty}(X,\theta;\phi)$), then so is $\varphi\land \psi$.
\end{proposition}
\begin{proof}
    The case of $\mathcal{E}^{\infty}(X,\theta;\phi)$ is trivial.

    We consider the case $\mathcal{E}(X,\theta;\phi)$. It follows from \cref{prop:landfinitecond1} that $\varphi\land \psi\in \PSH(X,\theta)$. By \cref{thm:diamond}, we have
    \[
        \int_X \theta_{\varphi\land\psi}^n\geq \int_X \theta_{\phi}^n.
    \]
    By the monotonicity theorem \cref{thm:mono}, equality holds. By \cref{thm:Pvarphidiffdef}, we conclude that
    \[
        P_{\theta}[\varphi\land \psi]=\phi.
    \]

    Finally, assume that $\varphi,\psi\in \mathcal{E}^1(X,\theta;\phi)$.
    We may assume that $\varphi,\psi\leq \phi$. For each $j\geq 1$, consider the approximations:
	\[
	\varphi_j\coloneqq \varphi\lor(\phi-j)\,,\quad \psi_j\coloneqq \psi\lor(\phi-j).
	\]
	By \cref{lma:rooftop_MAequation}, we can take $\gamma_j\in \mathcal{E}^{\infty}(X,\theta;\phi)$ solving the following equation:
	\[
	\theta_{\gamma_j}^n=\mathrm{e}^{\gamma_j-\varphi_j}\,\theta_{\varphi_j}^n+\mathrm{e}^{\gamma_j-\psi_j}\,\theta_{\psi_j}^n.
	\]
    In particular,
    \[
        \theta_{\gamma_j}^n\geq \mathrm{e}^{\gamma_j-\varphi_j}\,\theta_{\varphi_j}^n.
    \]
   We claim that this implies $\gamma_j\leq \varphi_j$. Indeed, by \cref{prop:compa_princ},
    \[
        \int_{\{\varphi_j<\gamma_j\}}\theta_{\varphi_j}^n
        \geq
        \int_{\{\varphi_j<\gamma_j\}}\theta_{\gamma_j}^n
        \geq
        \int_{\{\varphi_j<\gamma_j\}}\mathrm{e}^{\gamma_j-\varphi_j}\,\theta_{\varphi_j}^n.
    \]
    Hence
    \[
        \int_{\{\varphi_j<\gamma_j\}}\left(\mathrm{e}^{\gamma_j-\varphi_j}-1\right)\theta_{\varphi_j}^n=0.
    \]
    It follows that $\int_{\{\varphi_j<\gamma_j\}}\theta_{\varphi_j}^n=0$, and the domination principle \cref{thm:dom_prin_1} gives $\gamma_j\leq \varphi_j$. Similarly, $\gamma_j\leq \psi_j$. In particular, $\gamma_j\leq \varphi_j\land \psi_j$.
	We claim that
	\begin{equation}\label{eq:intgammajphi_unif_lower}
		\int_X (\gamma_j-\phi)\,\theta_{\gamma_j}^n >-C
	\end{equation}
    for some $C$ independent of $j$.

	Assume the claim is true for now.
	We get immediately that
    \[
    \sup_{X}\gamma_j=\sup_{X\setminus \{\phi=-\infty\}}(\gamma_j-\phi)\geq -\frac{C}{\int_X \theta_{\phi}^n},
    \]
    where the first equality follows from \cref{prop:model_sup_prop}.
    By \cref{prop:L1compa}, after possibly extracting a subsequence, we may assume that $\gamma_j\to \gamma\in \PSH(X,\theta)$ in $L^1$-topology.
    Then $\gamma\in \mathcal{E}^1(X,\theta;\phi)$ by \cref{lma:stab} and \cref{prop:E1charac}. Since $\gamma\leq\varphi\land\psi\leq\phi$, the monotonicity theorem \cref{thm:mono} gives
    \[
        \int_X\theta_{\varphi\land\psi}^n=\int_X\theta_\phi^n.
    \]
    Moreover, by the monotonicity of $E^\phi_\theta$,
    \[
        E^\phi_\theta(\varphi\land\psi)\geq E^\phi_\theta(\gamma)>-\infty.
    \]
    Thus $\varphi\land\psi\in\mathcal E^1(X,\theta;\phi)$ by \cref{prop:E1charac}.


	Now we prove the claim \eqref{eq:intgammajphi_unif_lower}.
	By symmetry, it suffices to prove
	\[
	\int_X (\phi-\gamma_j) \mathrm{e}^{\gamma_j-\varphi_j}\,\theta_{\varphi_j}^n\leq C.
	\]
	But note that
	\[
	\int_X (\phi-\gamma_j) \mathrm{e}^{\gamma_j-\varphi_j}\,\theta_{\varphi_j}^n=\int_X (\phi-\varphi_j) \mathrm{e}^{\gamma_j-\varphi_j}\,\theta_{\varphi_j}^n+\int_X (\varphi_j-\gamma_j) \mathrm{e}^{\gamma_j-\varphi_j}\,\theta_{\varphi_j}^n.
	\]
	But $x \mathrm{e}^{-x}\leq \mathrm{e}^{-1}$ for $x\geq 0$ (with maximum at $x=1$), so the second term is bounded; it remains to prove
	\[
	\int_X (\phi-\varphi_j) \mathrm{e}^{\gamma_j-\varphi_j}\,\theta_{\varphi_j}^n\leq C.
	\]
	As $\gamma_j\leq \varphi_j$, it suffices to prove
	\begin{equation}\label{eq:inproof2}
		\int_X (\phi-\varphi_j)\,\theta_{\varphi_j}^n\leq C.
	\end{equation}
    We compute, using \eqref{eq:Edefbdd} and the fact that the $k<n$ terms
    $\int_X(\varphi_j-\phi)\theta_{\varphi_j}^k\wedge\theta_\phi^{n-k}$ are all non-positive (since $\varphi_j\leq \phi$),
    \[
        \int_X (\varphi_j-\phi)\,\theta_{\varphi_j}^n\geq (n+1)E_{\theta}^{\phi}(\varphi_j).
    \]
	Since $\varphi_j\searrow\varphi$ and $E_\theta^{\phi}$ is continuous along decreasing sequences (\cref{prop:E_cont_decnet}), $E_\theta^{\phi}(\varphi_j)$ is bounded below uniformly in $j$. Thus, \eqref{eq:inproof2} follows.
\end{proof}



\begin{proposition}\label{prop:relativeEupperclosed}
    Let $\varphi,\psi\in \PSH(X,\theta)$ be potentials such that $\varphi\preceq \psi \preceq \phi$. Assume that $\varphi\in \mathcal{E}(X,\theta;\phi)$ (resp. $\mathcal{E}^1(X,\theta;\phi)$, $\mathcal{E}^{\infty}(X,\theta;\phi)$), then so is $\psi$.
\end{proposition}
\begin{proof}
    We may assume that $\varphi\leq \psi$.

    The case $\mathcal{E}^{\infty}(X,\theta;\phi)$ is trivial. The case $\mathcal{E}(X,\theta;\phi)$ follows from \cref{thm:mono}. The case $\mathcal{E}^1(X,\theta;\phi)$ follows from the characterization of $\mathcal{E}^1(X,\theta;\phi)$ in \cref{prop:E1charac} and the monotonicity of $E_{\theta}^{\phi}$ proved in \cref{lma:EmonoEinfty}.
\end{proof}

\begin{proposition}\label{prop:supsEE1}
    Let $(\varphi_i)_{i\in I}$ be a uniformly bounded from above non-empty family in $\mathcal{E}(X,\theta;\phi)$ (resp. $\mathcal{E}^1(X,\theta;\phi)$, $\mathcal{E}^{\infty}(X,\theta;\phi)$), then so is $\uscsup[i\in I]\varphi_i$.
\end{proposition}
\begin{proof}
Thanks to \cref{prop:relativeEupperclosed}, it suffices to show that
\[
\uscsup[i\in I]\varphi_i\preceq \phi.
\]
Since $\phi$ is model and $\varphi_i\preceq \phi$, we know that
\[
\varphi_i-\sup_X \varphi_i \leq \phi
\]
for any $i\in I$.
Since $(\varphi_i)_{i\in I}$ is uniformly bounded from above by assumption, our assertion follows.
\end{proof}




\begin{proposition}\label{prop:envrelfullmass}
    Let $\varphi,\psi\in \mathcal{E}(X,\theta;\phi)$. Then
    \[
        P_{\theta}[\varphi](\psi)=\uscsup[C > 0] (\varphi+C)\land \psi=\psi.
    \]
\end{proposition}
\begin{proof}
    The first equality follows from the definition of $P_{\theta}[\varphi](\psi)$. We focus on the second equality.

    Since for each $C\geq 0$,
    \[
        (\varphi\land \psi+C)\land \psi\leq (\varphi+C)\land \psi\leq \psi,
    \]
    we may replace $\varphi$ by $\varphi\land \psi$ (cf. \cref{prop:relrooftopclosed}) and assume that $\varphi\leq \psi$.

    Let
    \[
        \gamma\coloneqq \uscsup[C > 0] (\varphi+C)\land \psi.
    \]
    Observe that $\gamma\leq \psi$, and $\gamma\in \mathcal{E}(X,\theta;\phi)$ due to \cref{prop:relrooftopclosed} and \cref{prop:supsEE1}.
    Then \cref{cor:Pvarphisupport} guarantees that
    \[
    \int_{\{\gamma<\psi\} } \theta_{\gamma}^n=0.
    \]
    Therefore, we can apply the domination principle \cref{thm:dom_prin_1} to conclude that $\gamma=\psi$.
\end{proof}


\begin{lemma}\label{lma:Ediff}
    Let $\varphi,\psi\in \mathcal{E}^{\infty}(X,\theta;\phi)$. Define
    \[
    \eta_t\coloneqq \Bigl( (1-t)\varphi+t\psi \Bigr)\land \psi,\quad t\in [0,1].
    \]
    Then $E_{\theta}^{\phi}(\eta_t)$ is differentiable for $t\in [0,1]$ and
    \begin{equation}\label{eq:Iphit_diff}
    \frac{\mathrm{d}}{\mathrm{d}t}E_{\theta}^{\phi}(\eta_t)=\int_{X} \Bigl( \psi-\min\{\varphi,\psi\} \Bigr)\,\theta_{\eta_t}^n.
    \end{equation}
\end{lemma}
    The derivatives at $t=0$ or $1$ are understood as one-sided derivatives.
\begin{proof}
    Let us prove \eqref{eq:Iphit_diff} with right-derivative instead of derivative, and $t\in [0,1)$. The left-derivative case is completely parallel.

    For each $t\in [0,1]$, we let
    \[
    f_t=\min \{ (1-t)\varphi+t\psi, \psi\}.
    \]

    Fix $t\in [0,1)$ and $s>0$ small enough so that $t+s<1$. Thanks to \cref{prop:Ediffbdd}, we have
    \[
    \begin{aligned}
    E_{\theta}^{\phi}(\eta_{t+s})-E_{\theta}^{\phi}(\eta_{t})\leq & \int_X \left(\eta_{t+s}-\eta_t\right)\,\theta_{\eta_t}^n \\
    =& \int_X \left(\eta_{t+s}-f_t\right)\,\theta_{\eta_t}^n &\textup{by \cref{lma:rooftopMA}}\\
    \leq & \int_X \left(f_{t+s}-f_t \right)\,\theta_{\eta_t}^n\\
    =& s\int_X  \Bigl(\psi-\min\{\varphi,\psi\}\Bigr)\,\theta_{\eta_t}^n.
    \end{aligned}
    \]
    Similarly, we get
    \[
    E_{\theta}^{\phi}(\eta_{t+s})-E_{\theta}^{\phi}(\eta_{t})\geq s\int_X  \Bigl(\psi-\min\{\varphi,\psi\}\Bigr)\,\theta_{\eta_{t+s}}^n.
    \]
    Observe that $\eta_{t+s}$ converges in capacity (in fact, even uniformly) to $\eta_t$ as $s\to 0+$. The desired result \eqref{eq:Iphit_diff} is then just a consequence of \cref{thm:loc_conv_cap}.
\end{proof}



\begin{lemma}\label{lma:Eaffinediff}
    Let $\varphi,\psi\in \mathcal{E}^{1}(X,\theta;\phi)$, and set
    \[
        \varphi_t\coloneqq (1-t)\varphi+t\psi,\qquad t\in[0,1].
    \]
    Then the function
    \[
        t\mapsto E_{\theta}^{\phi}(\varphi_t)
    \]
    is differentiable on $[0,1]$, where the derivative at the endpoints is understood as the one-sided derivative. Moreover,
    \begin{equation}\label{eq:Eaffinediff}
        \frac{\mathrm d}{\mathrm dt}E_{\theta}^{\phi}(\varphi_t)
        =
        \int_X(\psi-\varphi)\,\theta_{\varphi_t}^n.
    \end{equation}
\end{lemma}
\begin{proof}
    By \cref{cor:E_classes_convex}, $\varphi_t\in \mathcal E^1(X,\theta;\phi)$ for all $t\in[0,1]$.
    For $a=0,\ldots,n$, put
    \[
        M_a=\int_X(\psi-\varphi)\,\theta_{\psi}^a\wedge\theta_{\varphi}^{n-a}.
    \]
    Note that each $M_a$ is finite by \cref{cor:mix_MA_energy_finite}. By the multi-linearity of the non-pluripolar product \itemref{itm:prop:npp1:6}, the quantity
    \[
        \int_X(\psi-\varphi)\,\theta_{\varphi_t}^n
    \]
    is therefore a polynomial in $t$ with coefficients given by the $M_a$'s. In particular, the right-hand side of \eqref{eq:Eaffinediff} is finite and continuous in $t$.

    We prove the formula for the right derivative at a point $t\in[0,1)$. Let $s>0$ be small enough so that $t+s\leq 1$. By the cocycle formula \cref{prop:cocycE1}, we have
    \[
    \begin{aligned}
        \frac{E_{\theta}^{\phi}(\varphi_{t+s})-E_{\theta}^{\phi}(\varphi_t)}{s}
        =&\frac{1}{(n+1)s}\sum_{j=0}^n
        \int_X(\varphi_{t+s}-\varphi_t)\,
        \theta_{\varphi_{t+s}}^j\wedge\theta_{\varphi_t}^{n-j} \\
        =&\frac{1}{n+1}\sum_{j=0}^n
        \int_X(\psi-\varphi)\,
        \theta_{\varphi_{t+s}}^j\wedge\theta_{\varphi_t}^{n-j}.
    \end{aligned}
    \]
    Again by \itemref{itm:prop:npp1:6}, for each fixed $j$ the last integral is a polynomial in $s$, expressed as a linear combination of the $M_a$'s with coefficients depending on $t$. Letting $s\to 0+$, it converges to
    \[
        \int_X(\psi-\varphi)\,\theta_{\varphi_t}^n.
    \]
    Therefore,
    \[
        \lim_{s\to0+}
        \frac{E_{\theta}^{\phi}(\varphi_{t+s})-E_{\theta}^{\phi}(\varphi_t)}{s}
        =
        \int_X(\psi-\varphi)\,\theta_{\varphi_t}^n.
    \]

    The proof for the left derivative is identical, using $s<0$ with $t+s\in[0,1]$. Thus the left and right derivatives agree for $t\in(0,1)$, and the endpoint derivatives are given by the same formula.
\end{proof}

In particular, we have a diamond-like inequality.
\begin{corollary}\label{cor:diam_E}
    Let $\varphi,\psi\in \mathcal{E}^1(X,\theta;\phi)$. Then
    \begin{equation}\label{eq:diam_E}
    E_{\theta}^{\phi}(\varphi\lor \psi)+E_{\theta}^{\phi}(\varphi\land \psi)\geq E_{\theta}^{\phi}(\psi)+E_{\theta}^{\phi}(\varphi).
    \end{equation}
\end{corollary}
\begin{proof}
    \textbf{Step~1}. We first reduce to the case where $\varphi,\psi \in \mathcal{E}^{\infty}(X,\theta;\phi)$.

    Assume that \eqref{eq:diam_E} is known in that case. For each $C>0$, we let
    \[
    \varphi_C=\varphi\lor (\phi-C),\quad \psi_C=\psi\lor (\phi-C).
    \]
    Then
    \[
    E_{\theta}^{\phi}(\varphi_C\lor \psi_C)+E_{\theta}^{\phi}(\varphi_C\land \psi_C)\geq E_{\theta}^{\phi}(\psi_C)+E_{\theta}^{\phi}(\varphi_C).
    \]
    As $C\to\infty$, we have
    \[
        \varphi_C\searrow\varphi,\quad \psi_C\searrow\psi,\quad \varphi_C\lor\psi_C\searrow\varphi\lor\psi,
    \]
    and, by the monotonicity and maximality of the rooftop,
    \[
        \varphi_C\land\psi_C\searrow\varphi\land\psi.
    \]
    Applying \cref{prop:E_cont_decnet} to these four decreasing sequences gives \eqref{eq:diam_E}.

    \textbf{Step~2}. We assume that $\varphi,\psi \in \mathcal{E}^{\infty}(X,\theta;\phi)$.

    Thanks to \cref{lma:Eaffinediff} and the plurifine locality of the non-pluripolar product \itemref{itm:prop:npp1:1}, we have
    \[
    \begin{aligned}
    E_{\theta}^{\phi}(\varphi\lor \psi)-E_{\theta}^{\phi}(\varphi)=&\int_0^1\int_X (\varphi\lor \psi-\varphi)\,\theta_{(1-t)\varphi+t(\varphi\lor \psi)}^n\,\mathrm{d}t\\
    =& \int_0^1 \int_{\{\psi>\varphi\}} (\psi-\varphi)\,\theta_{(1-t)\varphi+t\psi}^n\,\mathrm{d}t.
    \end{aligned}
    \]
    On the other hand,
    \[
    \begin{aligned}
        &E_{\theta}^{\phi}(\psi)-E_{\theta}^{\phi}(\varphi\land \psi)\\
        =& \int_0^1 \int_X (\psi-\min\{\varphi,\psi\})\,\theta_{((1-t)\varphi+t\psi)\land \psi}^n\,\mathrm{d}t\\
        \leq & \int_0^1 \int_{\{\varphi<\psi\}} (\psi-\varphi)\,\theta_{(1-t)\varphi+t\psi}^n\,\mathrm{d}t& \textup{by \cref{lma:rooftopMA}}.
    \end{aligned}
    \]
    Adding these equations together, we conclude \eqref{eq:diam_E}.
\end{proof}


\section{The \texorpdfstring{$\mathcal{I}$}{I}-envelope}\label{sec:Ienv}
From the algebraic point of view, a more natural envelope operator is given by the $\mathcal{I}$-envelope.

In this section, $X$ will denote a connected compact Kähler manifold of dimension $n$.


\subsection{\texorpdfstring{$\mathcal{I}$}{I}-equivalence}
\label{subsec:Iequiv}

\begin{proposition}\label{prop:Iequivchar}
    Given $\varphi,\psi\in \QPSH(X)$, the following are equivalent:
    \begin{enumerate}
        \item\label{itm:prop:Iequivchar:1} For any $k\in \mathbb{Z}_{>0}$, we have
        \[
            \mathcal{I}(k\varphi)=\mathcal{I}(k\psi);
        \]
        \item\label{itm:prop:Iequivchar:2} for any $\lambda\in \mathbb{R}_{>0}$, we have
            \[
            \mathcal{I}(\lambda\varphi)=\mathcal{I}(\lambda\psi);
            \]
        \item\label{itm:prop:Iequivchar:3} for any modification $\pi\colon Y\rightarrow X$ and any $y\in Y$, we have
            \[
                \nu(\pi^*\varphi,y)=\nu(\pi^*\psi,y);
            \]
        \item\label{itm:prop:Iequivchar:4} for any proper bimeromorphic morphism $\pi\colon Y\rightarrow X$ from a Kähler manifold and any $y\in Y$, we have
            \[
                \nu(\pi^*\varphi,y)=\nu(\pi^*\psi,y);
            \]
        \item\label{itm:prop:Iequivchar:5} for any prime divisor $E$ over $X$, we have
            \[
                \nu(\varphi,E)=\nu(\psi,E).
            \]
    \end{enumerate}
\end{proposition}
See \cref{def:primdiv} for the definition of prime divisors over $X$. Recall that in the whole book, a \emph{modification} of a compact complex space means a finite composition of blow-ups with smooth centers. This terminology is highly non-standard.

For simplicity, we usually write
\[
\nu(\varphi,y)=\nu(\pi^*\varphi,y)
\]
when $\pi$ is obvious from the context.

\begin{proof}
    \ref{itm:prop:Iequivchar:4} $\iff$ \ref{itm:prop:Iequivchar:5}. This follows from \cref{lma:blowupLelong}.

    \ref{itm:prop:Iequivchar:3} $\iff$ \ref{itm:prop:Iequivchar:5}. This follows from \cref{cor:primerealization}.

    \ref{itm:prop:Iequivchar:1} $\implies$ \ref{itm:prop:Iequivchar:5}. This follows from \cref{prop:Lelongnumfrommis}.

    \ref{itm:prop:Iequivchar:5} $\implies$ \ref{itm:prop:Iequivchar:2}. This follows from \cref{thm:valuativemulti}.

    \ref{itm:prop:Iequivchar:2} $\implies$ \ref{itm:prop:Iequivchar:1}. This is the restriction to positive integral values of $\lambda$.
\end{proof}


\begin{definition}\label{def:Iequiv}
    Given $\varphi,\psi\in \QPSH(X)$, we say they are \emph{$\mathcal{I}$-equivalent}\index{I-equivalent@$\mathcal{I}$-equivalence} and write $\varphi\sim_{\mathcal{I}}\psi$\index{$\sim_{\mathcal{I}}$} if the equivalent conditions in \cref{prop:Iequivchar} are satisfied.
\end{definition}
Clearly, $\sim_{\mathcal{I}}$ is an equivalence relation on $\QPSH(X)$.

\begin{proposition}\label{prop:Ienvbimero}
    Let $\pi\colon Y\rightarrow X$ be a proper bimeromorphic morphism from a Kähler manifold $Y$ to $X$. For $\varphi,\psi\in \QPSH(X)$, the following are equivalent:
    \begin{enumerate}
        \item\label{itm:prop:Ienvbimero:1} $\varphi\sim_{\mathcal{I}}\psi$;
        \item\label{itm:prop:Ienvbimero:2} $\pi^*\varphi\sim_{\mathcal{I}}\pi^*\psi$.
    \end{enumerate}
\end{proposition}
\begin{proof}
    \ref{itm:prop:Ienvbimero:1} $\implies$ \ref{itm:prop:Ienvbimero:2}. This follows from \itemref{itm:prop:Iequivchar:4}.

    \ref{itm:prop:Ienvbimero:2} $\implies$ \ref{itm:prop:Ienvbimero:1}. This follows from the simple fact that
    \[
        \mathcal{I}(k\varphi)=\pi_* \left( \omega_{Y/X}\otimes \mathcal{I}(k\pi^*\varphi)  \right),\quad \mathcal{I}(k\psi)=\pi_* \left( \omega_{Y/X}\otimes \mathcal{I}(k\pi^*\psi)  \right)
    \]
    for any $k\in \mathbb{Z}_{>0}$, where $\omega_{Y/X}$ is the relative dualizing sheaf. See \cite[Proposition~5.8]{Dem12}.
\end{proof}

\begin{proposition}\label{prop:Iequivmax}
    Let $\varphi,\varphi',\psi,\psi'\in \QPSH(X)$ and $\lambda>0$. Assume that $\varphi\sim_{\mathcal{I}}\psi$ and $\varphi'\sim_{\mathcal{I}}\psi'$. Then
    \[
        \varphi\lor \varphi'\sim_{\mathcal{I}} \psi\lor\psi',\quad \varphi+\varphi'\sim_{\mathcal{I}} \psi+\psi',\quad \lambda\varphi\sim_{\mathcal{I}}\lambda\psi.
    \]
    Similarly, if $(\varphi_i)_{i\in I}$, $(\psi_i)_{i\in I}$ are two non-empty uniformly bounded from above families in $\PSH(X,\theta)$ for some closed smooth real $(1,1)$-form $\theta$ on $X$ such that $\varphi_i\sim_{\mathcal{I}}\psi_i$ for all $i\in I$, then
    \[
    \uscsup[i\in I] \varphi_i\sim_{\mathcal{I}}\uscsup[i\in I] \psi_i.
    \]
\end{proposition}
\begin{proof}
    This follows from \cref{prop:Lelongmax} and \cref{cor:supsLelong}.
\end{proof}

\subsection{The definition of the \texorpdfstring{$\mathcal{I}$}{I}-envelope}
\label{subsec:Ienv_def}
Fix a smooth closed real $(1,1)$-form $\theta$ on $X$.
\begin{definition}\label{def:Ienv}
    Given $\varphi\in \PSH(X,\theta)$, we define its \emph{$\mathcal{I}$-envelope}\index{envelope!I-envelope@$\mathcal{I}$-envelope} as follows:
    \begin{equation}\label{eq:Ienvelopedef}
        P_{\theta}[\varphi]_{\mathcal{I}}\coloneqq \uscsup\{\psi\in \PSH(X,\theta): \psi\leq 0,\psi\sim_{\mathcal{I}} \varphi\}.
    \end{equation}
    \index{$P_{\theta}[\varphi]_{\mathcal{I}}$}
    If $\varphi=P_{\theta}[\varphi]_{\mathcal{I}}$, we say $\varphi$ is an \emph{$\mathcal{I}$-model potential}\index{potential!I-model potential@$\mathcal{I}$-model potential} (in $\PSH(X,\theta)$).
\end{definition}
Note that by \cref{prop:pshfunction_closedseq}, $P_{\theta}[\varphi]_{\mathcal{I}}\in \PSH(X,\theta)$.

\begin{proposition}\label{prop:Ienvindeptheta}
    Let $\theta'=\theta+\ddc g$ for some $g\in C^{\infty}(X)$. Then for any $\varphi\in \PSH(X,\theta)$, we have $\varphi'\coloneqq \varphi-g\in \PSH(X,\theta')$ and
    \[
        P_{\theta}[\varphi]_{\mathcal{I}}\sim P_{\theta'}[\varphi']_{\mathcal{I}}.
    \]
\end{proposition}
\begin{proof}
    By symmetry, it suffices to show that
    \[
        P_{\theta}[\varphi]_{\mathcal{I}}\preceq P_{\theta'}[\varphi']_{\mathcal{I}}.
    \]
    After replacing $g$ by $g+C$ for a suitable constant $C$, which only changes $\varphi'$ by an additive constant, we may assume that $g\geq 0$.
    Let $\psi\in \PSH(X,\theta)$ be such that $\psi\leq 0$ and $\psi\sim_{\mathcal{I}}\varphi$. Then $\psi'\coloneqq \psi-g$ lies in $\PSH(X,\theta')$, satisfies $\psi'\leq 0$, and is $\mathcal{I}$-equivalent to $\varphi'$ since adding a smooth function does not change multiplier ideals. Hence
    \[
        \psi-g=\psi'\leq P_{\theta'}[\varphi']_{\mathcal{I}}.
    \]
    It follows that
    \[
        \psi-\sup_X g\leq P_{\theta'}[\varphi']_{\mathcal{I}}.
    \]
    Taking the upper envelope over all such $\psi$ gives
    \[
        P_{\theta}[\varphi]_{\mathcal{I}}-\sup_X g\leq P_{\theta'}[\varphi']_{\mathcal{I}},
    \]
    which proves the desired comparison.
\end{proof}

\begin{proposition}\label{prop:Ienvelopebimero}
    Let $\pi\colon Y\rightarrow X$ be a proper bimeromorphic morphism from a Kähler manifold $Y$ to $X$. Then for $\varphi \in \PSH(X,\theta)$, we have
    \[
        P_{\pi^*\theta}[\pi^*\varphi]_{\mathcal{I}}=\pi^* P_{\theta}[\varphi]_{\mathcal{I}}.
    \]
\end{proposition}
\begin{proof}
    Let
    \[
        \eta=P_{\pi^*\theta}[\pi^*\varphi]_{\mathcal{I}}.
    \]
    First take a candidate $\psi\in \PSH(X,\theta)$ in \eqref{eq:Ienvelopedef}, so that $\psi\leq 0$ and $\psi\sim_{\mathcal{I}}\varphi$. Then $\pi^*\psi\leq 0$ and $\pi^*\psi\sim_{\mathcal{I}}\pi^*\varphi$ by \cref{prop:Ienvbimero}; hence $\pi^*\psi\leq \eta$. Taking the upper envelope over all such $\psi$ gives
    \[
        \pi^* P_{\theta}[\varphi]_{\mathcal{I}}\leq \eta.
    \]

    Conversely, let $\chi\in \PSH(Y,\pi^*\theta)$ be a candidate in the definition of $\eta$, so that $\chi\leq 0$ and $\chi\sim_{\mathcal{I}}\pi^*\varphi$. By \cref{prop:PSHpullbij}, there is a unique $\psi\in \PSH(X,\theta)$ such that $\chi=\pi^*\psi$. Since $\pi$ is surjective, $\psi\leq 0$. By \cref{prop:Ienvbimero}, the relation $\chi\sim_{\mathcal{I}}\pi^*\varphi$ implies $\psi\sim_{\mathcal{I}}\varphi$. Thus $\psi\leq P_{\theta}[\varphi]_{\mathcal{I}}$, and hence
    \[
        \chi\leq \pi^* P_{\theta}[\varphi]_{\mathcal{I}}.
    \]
    Taking the upper envelope over all such $\chi$ gives the reverse inequality.
\end{proof}

\begin{proposition}\label{prop:Ienvprojection}
    Let $\varphi\in \PSH(X,\theta)$. Then
    \[
        \varphi\sim_{\mathcal{I}}  P_{\theta}[\varphi]_{\mathcal{I}}.
    \]
    In particular,
    \[
        P_{\theta}\Bigl[ P_{\theta}[\varphi]_{\mathcal{I}} \Bigr]_{\mathcal{I}}=P_{\theta}[\varphi]_{\mathcal{I}}
    \]
    and the upper semicontinuous regularization in \eqref{eq:Ienvelopedef} is not necessary.
\end{proposition}
\begin{proof}
    In view of \cref{prop:Iequivchar}, it suffices to show that for $k \in \mathbb{Z}_{>0}$, we have
    \begin{equation}\label{eq:IenvelopepreservLelong}
        \mathcal{I}(k\varphi)=\mathcal{I}(kP_{\theta}[\varphi]_{\mathcal{I}}).
    \end{equation}

    By \cref{prop:Choquet}, we can find $\psi_i\in \PSH(X,\theta)$ ($i\in \mathbb{Z}_{>0}$) such that $\psi_i\leq 0$, $\psi_i\sim_{\mathcal{I}}\varphi$ for all $i\geq 1$ and
    \[
        \uscsup[i>0]\psi_i=P_{\theta}[\varphi]_{\mathcal{I}}.
    \]
    By \cref{prop:Iequivmax}, we may replace $\psi_i$ by $\psi_1\lor\cdots\lor\psi_i$ and assume that the sequence $\psi_i$ is increasing. In this case, it follows from the strong openness theorem \cref{thm:strongopen} that for each $k\in \mathbb{Z}_{>0}$, we have
    \[
        \mathcal{I}(k\varphi)=\mathcal{I}(k\psi_j)=\mathcal{I} (kP_{\theta}[\varphi]_{\mathcal{I}})
    \]
    for $j$ large enough.
\end{proof}




\begin{definition}\label{def:volqpsh}
    Let $\varphi\in \PSH(X,\theta)$, we define the \emph{volume}\index{volume}\footnote{We choose to call this quantity the \emph{volume} instead of the \emph{$\mathcal{I}$-volume} so that the terminology is consistent with the line bundle case.} $\vol(\theta,\varphi)$ as
    \[
        \vol(\theta,\varphi)=\int_X \Bigl(\theta+\ddc P_{\theta}[\varphi]_{\mathcal{I}}  \Bigr)^n.
    \]
    \index{$\vol(\theta,\varphi)$}
\end{definition}

\begin{proposition}\label{prop:voldeponlyoncurr}
    Let $\theta'=\theta+\ddc g$ for some $g\in C^{\infty}(X)$. Then for any $\varphi\in \PSH(X,\theta)$, we have $\varphi'=\varphi-g\in \PSH(X,\theta')$ and
    \[
        \vol(\theta,\varphi)=  \vol\left(\theta',\varphi'\right).
    \]
\end{proposition}
\begin{proof}
    By \cref{prop:Ienvindeptheta}, the currents $\theta+\ddc P_\theta[\varphi]_{\mathcal I}$ and $\theta'+\ddc P_{\theta'}[\varphi']_{\mathcal I}$ have the same singularity type in the same cohomology class. Applying the monotonicity theorem \cref{thm:mono} gives equality of their non-pluripolar masses.
\end{proof}

In view of \cref{prop:voldeponlyoncurr}, the volume $\vol(\theta,\varphi)$ depends only on the current $\theta_{\varphi}$, and we may write
\begin{equation}\label{eq:volcurrdef}
    \vol \theta_{\varphi}=\vol(\theta,\varphi).
\end{equation}
\index{$\vol \theta_{\varphi}$}

\begin{definition}\label{def:volcohclass}
    Let $\alpha\in \mathrm{H}^{1,1}(X,\mathbb{R})$ be a pseudo-effective class. The \emph{volume}\index{volume!volume of a pseudo-effective class} $\vol \alpha$ of $\alpha$ is defined as
    \[
    \vol \alpha=\vol T_{\min},
    \]
    where $T_{\min}$ is a current with minimal singularities in $\alpha$. Note that $\vol \alpha$ is independent of the choice of $T_{\min}$, thanks to \cref{thm:mono}.
\end{definition}

Let us recall the following elementary result for later use.
\begin{proposition}\label{prop:cones_modif2}
    Let $\pi\colon Y\rightarrow X$ be a proper bimeromorphic morphism from a Kähler manifold $Y$.
    \begin{enumerate}
        \item\label{itm:prop:cones_modif2:1} For any big class $\alpha\in \mathrm{H}^{1,1}(X,\mathbb{R})$, $\pi^*\alpha$ is big. Moreover, $\vol \pi^*\alpha=\vol \alpha$.
        \item\label{itm:prop:cones_modif2:2} For any big class $\beta\in \mathrm{H}^{1,1}(Y,\mathbb{R})$, $\pi_*\beta$ is big. Moreover, $\vol \pi_*\beta\geq \vol \beta$.
    \end{enumerate}
\end{proposition}
\begin{proof}
    \ref{itm:prop:cones_modif2:1} Take a current $T_{\min}$ with minimal singularities in $\alpha$. Then $\pi^*T_{\min}$ is a current with minimal singularities in $\pi^*\alpha$. Our assertion follows from \cref{prop:nppmassinv}.

    \ref{itm:prop:cones_modif2:2} Let $S_{\min}$ be a current with minimal singularities in $\beta$. Then $\pi_*S_{\min}$ is a positive current in the cohomology class $\pi_*\beta$. A current with minimal singularities in $\pi_*\beta$ is less singular than $\pi_*S_{\min}$, hence by the monotonicity theorem \cref{thm:mono}
    \[
        \vol(\pi_*\beta)\geq \int_X\left(\pi_*S_{\min}\right)^n.
    \]
    By \cref{prop:nppmassinv},
    \[
        \int_X \left(\pi_*S_{\min}\right)^n\geq \int_Y S_{\min}^n=\vol\beta.
    \]
    We conclude $\vol \pi_*\beta\geq \vol \beta$.
\end{proof}


The $\mathcal{I}$-envelope and the $P$-envelope are related in a simple manner.
\begin{proposition}\label{prop:PandPI}
    Let $\varphi\in \PSH(X,\theta)$. Then
    \[
        P_{\theta}[\varphi]\leq P_{\theta}[\varphi]_{\mathcal{I}},\quad \varphi \sim_{\mathcal{I}}P_{\theta}[\varphi].
    \]
\end{proposition}
As a consequence,
\begin{equation}\label{eq:PthetaIP_comp}
    P_{\theta}\Bigl[P_{\theta}[\varphi]\Bigr]_{\mathcal{I}}=P_{\theta}[\varphi]_{\mathcal{I}}.
\end{equation}
\begin{proof}
    It suffices to show that $\varphi\sim_{\mathcal{I}} P_{\theta}[\varphi]$. Namely, for each $k\in \mathbb{Z}_{>0}$, we have
    \begin{equation}\label{eq:IkvarphiIkP}
        \mathcal{I}(k\varphi)=\mathcal{I}\Bigl(kP_{\theta}[\varphi]\Bigr).
    \end{equation}
    Fix $k$ for now.
    It follows from \eqref{eq:Penvsups} and the strong openness theorem \cref{thm:strongopen} that
    \[
        \mathcal{I}\Bigl(kP_{\theta}[\varphi]\Bigr)=\mathcal{I}\Bigl((k\varphi+C)\land kV_{\theta}\Bigr),
    \]
    when $C$ is large enough. Since $(k\varphi+C)\land kV_{\theta}\sim k\varphi$, we have
    \[
        \mathcal{I}\Bigl((k\varphi+C)\land kV_{\theta}\Bigr)=\mathcal{I}(k\varphi)
    \]
    and \eqref{eq:IkvarphiIkP} follows.
\end{proof}

In particular, we obtain an interesting relation between the non-pluripolar mass and the volume.
\begin{corollary}\label{cor:compnppmassandvol}
    Let $\varphi\in \PSH(X,\theta)$. Then
    \[
        \int_X \theta_{\varphi}^n\leq \vol \theta_{\varphi}.
    \]
\end{corollary}
The reverse inequality fails in general, see \cref{ex:BBJ}.
\begin{proof}
    This follows from \cref{prop:PandPI}, \cref{thm:mono} and \cref{prop:Ppresmass}.
\end{proof}


We note the following special case:
\begin{proposition}\label{prop:analysingcompPandPI}
    Let $\varphi\in \PSH(X,\theta)$. Assume that $\varphi$ has analytic singularities. Then
    \begin{equation}\label{eq:ana_sing_phi_PI_eq}
        \varphi\sim P_{\theta}[\varphi] \sim P_{\theta}[\varphi]_{\mathcal{I}}.
    \end{equation}
    Hence, $P_{\theta}[\varphi]$ and $P_{\theta}[\varphi]_{\mathcal{I}}$ both have analytic singularities.

    In particular,
    \begin{equation}\label{eq:ana_sing_mass_vol}
        \int_X \theta_{\varphi}^n= \vol \theta_{\varphi}.
    \end{equation}
\end{proposition}
\begin{proof}
    First observe that \eqref{eq:ana_sing_mass_vol} follows from \eqref{eq:ana_sing_phi_PI_eq} and \cref{thm:mono}. It remains to establish \eqref{eq:ana_sing_phi_PI_eq}.

    In view of \cref{prop:PandPI}, it suffices to show that
    \begin{equation}\label{eq:Pprecvarphitemp1}
        P_{\theta}[\varphi]_{\mathcal{I}}\preceq \varphi.
    \end{equation}
    By \cref{prop:Ienvelopebimero,prop:Penvbimero} and \cref{thm:resolvelogsing}, after passing to a log resolution we may assume that $\varphi$ has log singularities along an effective $\mathbb Q$-divisor $D$. This reduction does not change the assertion, since both the $P$- and the $\mathcal I$-envelope are compatible with bimeromorphic pull-back.

    After multiplying $\theta$ and $\varphi$ by a sufficiently divisible positive integer, and then dividing the final comparison by the same integer, we may assume that $D$ is integral. In this case the Bergman-type quasi-equisingular approximants of $\varphi$ associated to the ideals $\mathcal I(2^j\varphi)$ have the same singularity type as $\varphi$.

    Since
    \[
        P_\theta[\varphi]_{\mathcal I}\sim_{\mathcal I}\varphi
    \]
    by \cref{prop:Ienvprojection}, we have
    \[
        \mathcal I(2^jP_\theta[\varphi]_{\mathcal I})=\mathcal I(2^j\varphi).
    \]
    Thus the quasi-equisingular approximants $\eta_j$ of $P_\theta[\varphi]_{\mathcal I}$ and $\varphi_j$ of $\varphi$ may be chosen with the same singularity type $(2^{-j},\mathcal I(2^j\varphi))$.
    Consequently $\eta_j\sim\varphi_j\sim\varphi$ for $j\gg1$. Since $\eta_j\searrow P_\theta[\varphi]_{\mathcal I}$, choosing one such large $j$ gives
    \[
        P_\theta[\varphi]_{\mathcal I}\leq \eta_j\leq \varphi+C_j,
    \]
    and hence \eqref{eq:Pprecvarphitemp1} holds.
\end{proof}



\subsection{Properties of the \texorpdfstring{$\mathcal{I}$}{I}-envelope}
\label{subsec:Ienv_prop}
Let $\theta,\theta_1,\theta_2$ be smooth closed real $(1,1)$-forms on $X$.

We have the following properties of the $\mathcal{I}$-envelope.
\begin{proposition}\label{prop:PIconc}
    \leavevmode
    \begin{enumerate}
        \item\label{itm:prop:PIconc:1} Suppose that $\varphi\in \PSH(X,\theta)$ and $\lambda\in \mathbb{R}_{>0}$. Then
            \[
                P_{\lambda\theta}[\lambda\varphi]_{\mathcal{I}}=\lambda P_{\theta}[\varphi]_{\mathcal{I}}.
            \]
        \item\label{itm:prop:PIconc:2} Suppose that $\varphi_1\in \PSH(X,\theta_1)$ and $\varphi_2\in \PSH(X,\theta_2)$. Then
            \[
                P_{\theta_1+\theta_2}[\varphi_1+\varphi_2]_{\mathcal{I}}\geq P_{\theta_1}[\varphi_1]_{\mathcal{I}}+P_{\theta_2}[\varphi_2]_{\mathcal{I}}.
            \]
        \item\label{itm:prop:PIconc:3}
        Suppose that $\varphi_1\in \PSH(X,\theta_1)$ and $\varphi_2\in \PSH(X,\theta_2)$. Then
            \[
                P_{\theta_1+\theta_2}[\varphi_1+\varphi_2]_{\mathcal{I}}\sim_{\mathcal{I}} P_{\theta_1}[\varphi_1]_{\mathcal{I}}+P_{\theta_2}[\varphi_2]_{\mathcal{I}}.
            \]
        \item\label{itm:prop:PIconc:4}
            Suppose that $\varphi_1,\varphi_2\in \PSH(X,\theta)$. Then
                \[
                    P_{\theta}[\varphi_1\lor \varphi_2]_{\mathcal{I}}\sim_{\mathcal{I}} P_{\theta}[\varphi_1]_{\mathcal{I}}\lor P_{\theta}[\varphi_2]_{\mathcal{I}}.
                \]
    \end{enumerate}
\end{proposition}
\begin{proof}
    \ref{itm:prop:PIconc:1} This follows by rescaling the candidates in the definition of the $\mathcal{I}$-envelope.

    \ref{itm:prop:PIconc:2} Suppose that $\psi_1\in \PSH(X,\theta_1)$ and $\psi_2\in \PSH(X,\theta_2)$ satisfy
    \[
        \psi_i\leq 0,\quad \psi_i \sim_{\mathcal{I}} \varphi_i
    \]
    for $i=1,2$. Then thanks to \cref{prop:Iequivmax},
    \[
        \psi_1+\psi_2\leq 0,\quad \psi_1+\psi_2 \sim_{\mathcal{I}} \varphi_1+\varphi_2.
    \]
    It follows that
    \[
        \psi_1+\psi_2\leq P_{\theta_1+\theta_2}[\varphi_1+\varphi_2]_{\mathcal{I}}.
    \]
    Since $\psi_1$ and $\psi_2$ are arbitrary, we conclude.

    \ref{itm:prop:PIconc:3} and \ref{itm:prop:PIconc:4}: These follow easily from \cref{prop:Ienvprojection} and \cref{prop:Iequivmax}.
\end{proof}

\begin{lemma}\label{lma:PIenvmono1}
    Let $\varphi,\psi\in \PSH(X,\theta)$. Assume that $\varphi\preceq \psi$. Then
    \[
        P_{\theta}[\varphi]_{\mathcal{I}}\leq P_{\theta}[\psi]_{\mathcal{I}}.
    \]
\end{lemma}
\begin{proof}
    We may assume that $\psi\leq 0$.
    Observe that $P_{\theta}[\varphi]_{\mathcal{I}}\lor \psi\sim_{\mathcal{I}} \psi$ as a consequence of \cref{prop:Lelongmax} and \cref{prop:Ienvprojection}. So
    \[
    P_{\theta}[\varphi]_{\mathcal{I}}\lor \psi\leq P_{\theta}[\psi]_{\mathcal{I}}.
    \]
    The desired inequality follows.
\end{proof}

\begin{proposition}\label{prop:decnetmodelPI}
    Consider a decreasing net $(\varphi_i)_{i\in I}$ of model potentials, with $\varphi_i\in \PSH(X,\theta)_{>0}$ for all $i\in I$. Set
    \[
    \varphi\coloneqq \inf_{i\in I}\varphi_i.
    \]
    Suppose that $\varphi\not\equiv -\infty$ and $\int_X \theta_{\varphi}^n>0$. Then
    \[
        \inf_{i\in I} P_{\theta}[\varphi_i]_{\mathcal{I}}=P_{\theta}[\varphi]_{\mathcal{I}}.
    \]
    In particular, if the $\varphi_i$'s are all $\mathcal{I}$-model, then so is $\varphi$.
\end{proposition}
\begin{proof}
    Let $\eta=\inf_{i\in I} P_{\theta}[\varphi_i]_{\mathcal{I}}$. We have $\eta\geq P_{\theta}[\varphi]_{\mathcal{I}}$ as a consequence of \cref{lma:PIenvmono1}. In particular, $\eta\in \PSH(X,\theta)$.

    By \cref{prop:vol_limit_model}, we have
    \[
        \lim_{i\in I} \int_X\theta_{\varphi_i}^n=\int_X\theta_{\varphi}^n>0.
    \]
    We may assume that for each $i\in I$, $\int_X\theta_{\varphi_i}^n>\int_X\theta_{\varphi}^n$, since otherwise $\varphi_i=\varphi$ by \cref{thm:Pvarphidiffdef}, and hence $P_{\theta}[\varphi_i]_{\mathcal{I}}=P_{\theta}[\varphi]_{\mathcal{I}}$ due to \eqref{eq:PthetaIP_comp}; the same holds for all $j\geq i$, and there is nothing to prove.

    So by \cref{lma:pathoenvelope}, we can find a decreasing net $\epsilon_i\searrow 0$ ($i\in I$) with $\epsilon_i\in (0,1)$ and $\psi_i\in \PSH(X,\theta)$ such that for all $i\in I$,
    \[
        (1-\epsilon_i) \varphi_i+\epsilon_i\psi_i\leq \varphi.
    \]
    We can, for example, take
    \[
        \epsilon_i^{-1}=\frac{1}{2}\left( 1+\left(\frac{\int_X \theta_{\varphi_i}^n}{\int_X \theta_{\varphi_i}^n-\int_X \theta_{\varphi}^n}\right)^{1/n} \right).
    \]

    By \cref{prop:PIconc} and \cref{lma:PIenvmono1}, we have
    \[
        \eta+\epsilon_i P_{\theta}[\psi_i]_{\mathcal{I}}\leq (1-\epsilon_i) \eta+\epsilon_i P_{\theta}[\psi_i]_{\mathcal{I}}\leq (1-\epsilon_i) P_{\theta}[\varphi_i]_{\mathcal{I}}+\epsilon_i P_{\theta}[\psi_i]_{\mathcal{I}}\leq P_{\theta}[\varphi]_{\mathcal{I}}.
    \]
    Each $P_{\theta}[\psi_i]_{\mathcal{I}}$ has supremum $0$, so \cref{prop:L1compa} shows that the family is bounded in $L^1$. Hence $\epsilon_iP_{\theta}[\psi_i]_{\mathcal{I}}\to 0$ in $L^1$, and the preceding inequality gives $\eta\leq P_{\theta}[\varphi]_{\mathcal{I}}$ almost everywhere. By \cref{prop:pshfuncdetdense}, the inequality holds everywhere.
\end{proof}

\begin{proposition}\label{prop:incnetmodelPI}
    Let $(\varphi_i)_{i\in I}$ be an increasing net in $\PSH(X,\theta)_{>0}$ uniformly bounded from above. Let $\varphi\coloneqq \uscsup[i\in I]\varphi_i$. Then
    \[
        \uscsup[i\in I] P_{\theta}[\varphi_i]_{\mathcal{I}}=P_{\theta}[\varphi]_{\mathcal{I}}.
    \]
    In particular, if the $\varphi_i$'s are all $\mathcal{I}$-model, then so is $\varphi$.
\end{proposition}
\begin{proof}
    Let $\eta=\uscsup[i\in I] P_{\theta}[\varphi_i]_{\mathcal{I}}$. Then $\eta\leq P_{\theta}[\varphi]_{\mathcal{I}}$ as a consequence of \cref{lma:PIenvmono1}.

    By \cref{cor:incseqnppcont}, we have
    \[
        \lim_{i\in I} \int_X\theta_{\varphi_i}^n=\int_X\theta_{\varphi}^n>0.
    \]
    We may assume that for each $i\in I$, $\int_X\theta_{\varphi_i}^n<\int_X\theta_{\varphi}^n$, since otherwise $P_{\theta}[\varphi_i]=P_{\theta}[\varphi]$ by \cref{thm:Pvarphidiffdef}, and hence $P_{\theta}[\varphi_i]_{\mathcal{I}}=P_{\theta}[\varphi]_{\mathcal{I}}$ due to \eqref{eq:PthetaIP_comp}; the same holds for all $j\geq i$, and there is nothing to prove.

    So by \cref{lma:pathoenvelope}, we can find a decreasing net $\epsilon_i\searrow 0$ ($i\in I$) as in \eqref{eq:epsilon_inetchoice} with $\epsilon_i\in (0,1)$ and $\psi_i\in \PSH(X,\theta)$ such that for all $i\in I$,
    \[
        (1-\epsilon_i) \varphi+\epsilon_i\psi_i\leq \varphi_i.
    \]
    By \cref{prop:PIconc} and \cref{lma:PIenvmono1}, we have
    \[
        P_{\theta}[\varphi]_{\mathcal{I}}+\epsilon_i P_{\theta}[\psi_i]_{\mathcal{I}}\leq (1-\epsilon_i) P_{\theta}[\varphi]_{\mathcal{I}}+\epsilon_i P_{\theta}[\psi_i]_{\mathcal{I}}\leq P_{\theta}[\varphi_i]_{\mathcal{I}}\leq \eta.
    \]
    As above, $\epsilon_iP_{\theta}[\psi_i]_{\mathcal{I}}\to 0$ in $L^1$. It follows that $P_{\theta}[\varphi]_{\mathcal{I}}\leq \eta$ almost everywhere, hence everywhere by \cref{prop:pshfuncdetdense}.
\end{proof}


\begin{remark}\label{rmk:envelopes:L2443}
    One could also define the following interpolation between the $\mathcal{I}$-envelope and the $P$-envelope: Suppose $\varphi\in \PSH(X,\theta)_{>0}$, $j\in \{0,\ldots,n\}$. Then we let
    \[
    \begin{split}
    P_{\theta,j}[\varphi]\coloneqq \uscsup \left\{ \psi\in \PSH(X,\theta): \psi\leq 0,\varphi\preceq \psi, \int_X \theta_{\varphi}^j\wedge \theta_{P_{\theta}[\varphi]_{\mathcal{I}}}^{n-j}\right.\\
    \left.=\int_X \theta_{\psi}^j\wedge \theta_{P_{\theta}[\psi]_{\mathcal{I}}}^{n-j} \right\}.
    \end{split}
    \]
    Based on the techniques developed in \cref{chap:comp}, one could show that $P_{\theta,j}[\bullet]$ is a projection operator. The cases $j=0$ and $j=n$ recover the $\mathcal{I}$- and $P$-envelopes respectively; the intermediate operators $P_{\theta,j}$ have not, to our knowledge, been systematically studied in the literature.
\end{remark}



%\bibliographystyle{sgptalpha}
\bibliography{ComplexGeometry}

